%%%%%%%%%%%%%%%%%%%%%definitions%%%%%%%%%%%%%%%%%%%%%%%%%%%%%%%%%%%%%%%

\input{../common/header.tex}
\input{../common/newcommands.tex}
\usepackage{minted}

%%%%%%%%%%%%%%%%%%%%%%%%%%%%%DOCUMENT%%%%%%%%%%%%%%%%%%%%%%%%%%%%%%%%%%%%%%%
\begin{document}

\title{
The full-F electromagnetic model in toroidal geometry \textsc{Feltor}}
\author{ M.~Wiesenberger and M.~Held}
\maketitle

\begin{abstract}
Thermal \textsc{Feltor} is a model for global 3d thermal electromagnetic full-F gyro-fluid simulations.

\noindent
This document describes the (usage of) programs
contained in \mintinline{bash}{path/to/feltor/src/thermal}
and connects
the implemented formulas to relevant journal publications.

The thermal code uses many common modules with the feltor code for which the same documentations apply. In such cases we refer the reader to the corresponding sections in the `feltor/src/feltor/feltor.tex` file which compiles to `feltor.pdf` for short.
\end{abstract}
\tableofcontents

\section{Quick-start}
The thermal code is compiled and used completely analogously to the feltor code so we refer to the Quick-start section in `feltor.pdf` file for compilation and usage. Just replace `feltor` with `thermal` everywhere.

However, there currently is no glfw build available.
\section{The magnetic field} \label{sec:magnetic}
The thermal code uses the same magnetic field as the feltor code so we refer to the `The magnetic field` section in `feltor.pdf`.


%%%%%%%%%%%%%%%%%%%%%%%%%%%%%%%%%%%%%%%%%%%%%%%%%%%%%%%%%%%%%%%%%%%%%%%%%%%%%%%%%
\section{The model} \label{sec:model}
For the model we refer to the overleaf repository for a detailed description.

\section{The Input file} \label{sec:input_file}
Input file format: \href{https://en.wikipedia.org/wiki/JSON}{json}

%%%%%%%%%%%%%%%%%%%%%%%%%%%%%%%%%%%%%
\subsection{Magnetic field} \label{sec:geometry_file}
The section `Magnetic field` in `feltor.pdf` applies verbatim except:

\begin{minted}[texcomments]{js}
"magnetic_field":
{
    "curvmode": "true"// NO LONGER VALID!
}
\end{minted}

%%%%%%%%%%%%%%%%%%%%%%%%%
\subsection{Spatial grid} \label{sec:spatial}
The section `Spatial grid` in `feltor.pdf` applies verbatim.
%%%%%%%%%%%%%%%%%%%%%%%%%%%%%%%%%%%%%%%%%%%%%%%%%%%%%
\subsection{The boundary region} \label{sec:boundary}
The section `The boundary region` in `feltor.pdf` applies verbatim.
The only exception is when the values for wall and sheath boundary conditions are applied.
The wall density is a list of values
\begin{minted}[texcomments]{js}
"boundary":
{
    "wall" :
    {
        "bc" : "fixed", // Dirichlet boundary
        "nwall" : [0.1,0.1], // $N_{s,w} >0$ (see also background and minne)
        // length must equal number of species (see \ref{sec:species})
        // Quasineutrality must be observed $\sum_s z_s n_{s,wall} = 0$
        "uwall" : 0.0, // $U_{\parallel,w}$
        "twall" : 1, // $T_w > 0$ wall temperature
        "qwall" : 0, // $Q_{w}$
    }
    "wall":
    {
        "bc" : "floating" // Neumann boundary
    }
    "sheath":
    {
        "type": "bohm" //INVALID, only "wall" or "insulating" are allowed
    }
}
\end{minted}
Due to collisions, there can only be one wall temperature for Dirichlet boundary condition.
Note the restriction
\begin{align}
    \sum_s z_s N_{s,w} = 0
\end{align}
and that $T_{w}>0$ and $N_{s,w}>0$ must be strictly positive.

For the floating wall we set
\begin{align}
    N_w =& \frac{\int \d V \chi_w (1-\chi_w) N }{ \int \d V \chi_w ( 1-\chi_w)} \\
    P_w =& \frac{\int \d V \chi_w (1-\chi_w) (P_\parallel + 2P_\perp) }{ 3\int \d V \chi_w ( 1-\chi_w) }
\end{align}

The second exception is the box boundary conditions section:
%%%%%%%%%%%%%%%%%%%%%%%%%%%%%%%%%%%%%%%%%
\subsubsection{Box boundary conditions}
\begin{tcolorbox}[title=Note]
The box boundary conditions are hardcoded and cannot be changed in the input file.
As long as the penalization method is used they are in fact irrelevant.
\end{tcolorbox}
For all quantities except $\phi$ and $A_\parallel$ we choose Neumann zero boundary
condition on the bounding box.
The boundary conditions for the potential and the magnetic potential are not
penalized and thus hold on the bounding box. We choose (hardcode)
\begin{align}
\phi = 0
\text{ and }  \hat n \cn A_{\parallel} = 0
\end{align}
where $\hat n$ is the normal vector to the boundary.

%%%%%%%%%%%%%%%%%%%%%%%%%%%%%%%
\subsection{The species block: Initial conditions and sources} \label{sec:species}
The species block replaces the "init" and "source" block in feltor.

In order to setup a number of species the input file needs to contain the "species" field, which
consists of a list of dictionaries describing the physical parameters, the initialisation and the sources
for each species in the simulation. The first species must be the electron species and its mass
    must be negligible ( $\mu <10^{-3}$)
\begin{minted}[texcomments]{js}
"species":
[   // list of dictionaries; its length determines the number of species
    {
        "name" : "e", // determines field names in output file
        "mu" : 2.72121e-4, // mass relative to reference mass: $m_s/m_i$.
        // results in the neglect of gyro-radius effects for that species (s.a. \ref{sec:physical})
        "z" : -1, //charge number of species (can be floating point)
        "kappa" : 3.2, // Braginskii numerical prefactor for heat flux ($\kappa_i=3.9$)
        "pi" : 0.73, // Braginksii numerical prefactor for viscosity ($\pi_i = 0.96$)
        "init" : {
            //... describe initialisation of fields for this species \ref{sec:init}
        },
        "source" : {
            // ... describe source terms for this species \ref{sec:sources}
        }
    },
    {
        // ... more species
    }
]
\end{minted}


\subsubsection{Initialisation} \label{sec:init}

See the file {\tt init.h} to add your own custom condition.

We have the possibility to
initialize a species with the output of a previously
run simulation. Note that this can only be done for all species, i.e. it is not possible
to initialise one species from a file and the rest with a profile say.
Also, all species names, mass charge, kappa and pi values in the input file must be identical to the ones in the restart file.
% We have to see how this is implemented to see if we can add more / less flexibiltiy here
% also I don't foresee this option to be used much
\begin{minted}[texcomments]{js}
"init":
{
    "type" : "restart",
    // only the "file" parameter is needed (also see Section \ref{sec:restart})
    // Whenever any one species is restarted, all species are restarted
    // The species names and parameters in the previous simulation must match the current ones.
    "file" : "path/to/file.nc"
    // the file to use as a restart for $N_s$, $T_{\parallel,s}$, $T_{\perp,s}$, $U_{\parallel,s}$, $Q_{\parallel,s}$ and $Q_{\perp,s}$
}
\end{minted}

In contrast, we can initialize the fields $N_s$, $P_{\parallel,s}$, $P_{\perp,s}$, $U_{\parallel,s}$, $Q_{\parallel,s}$ and $Q_{\perp,s}$
directly:
\begin{minted}[texcomments]{js}
"init":
{
    "type" : "fields",
    "density":
    {
        //...
    },
    "pperp":
    {
        //...
    },
    "ppara":
    {
        //...
    },
    "velocity":
    {
        //...
    },
    "qperp":
    {
        //...
    },
    "qpara"
    {
        //...
    }
}
\end{minted}

\noindent
The "density", "pperp" and "ppara" fields all allow the following:
\begin{minted}[texcomments]{js}
"init":
{
    "density" : // or "pperp" or "ppara"
    {
        "type" : "const"  // $N_s = const$
        "background" : 0.1, // $const>0$
        // same as choosing "const" profile with zero ntilde and no damping
        //(see also minne and nwall)
    }
    // OR
    "density" : // or "pperp" or "ppara"
    {
        "type" : "profile",  // $n_s = n_{bg} + (n_{prof}(R,Z) + \tilde n(R,Z,\varphi) - n_{bg}) D(R,Z)$
        "ntilde":
        {
            // see Ntilde section in `feltor.pdf`
        },
        "profile":
        {
            // see Profile section in `feltor.pdf`
        },
        "damping":
        {
            // see Damping section in `feltor.pdf`
        }
    }
}
\end{minted}
\begin{tcolorbox}[title=Note]
    Even though the parameter reads "ntilde" it applies equally to pressure fluctuations.
    Also we here initialise the physical, particle quantities rather than the gyro-centre quantities
\end{tcolorbox}
For the velocity and heat fluxes we have
\begin{minted}[texcomments]{js}
"init":
{
    "velocity" :
    {
        "type" : "zero",  // $U_{\parallel}(R,Z,\varphi,0) = 0$
    },
    "qperp":
    {
        "type" : "zero" // currently the only option for $Q_\perp$
    },
    "qpara":
    {
        "type" : "zero" // currently the only option for $Q_\parallel$
    }
}
\end{minted}
\noindent
For the velocity to fulfill sheath boundary conditions we initialize a linear profile between the divertor plates.
This is achieved by integrating fieldlines towards the positive and negative divertor.
\begin{minted}[texcomments]{js}
"init":
{
    "velocity" :
    {
        "type" : "linear_cs" // $ U_{\parallel,i}= s\sqrt{\frac{T_{i,\parallel} + z_i T_{e,\parallel}}{\mu_i}}$ in the SOL
        // where "s" is the fieldline coordinate defined in section\ref{sec:sheath}
        // if the sheath type is "none" then $s=0$
        // this type does not work for electrons
        // (electrons are the first species)
    }
}
\end{minted}
\noindent


\paragraph{Quasineutrality and insulating sheath}
We need to observe the quasineutrality
\begin{align}
\sum_s z_s n_s = 0
\end{align}
and for the sheath we need to have zero outflow at the sheath
\begin{align}
\sum_s z_s n_s u_{\parallel,s}|_{sh} = 0
\end{align}

In order to achieve this a single species (usually the electrons) can be automatically determined by the
above conditions setting
\begin{minted}[texcomments]{js}
"init":
{
    "density" :
    {
        "type" : "quasineutral"  // $n_s = -\sum_{i,i\neq s} z_i n_i/ z_s$
    },
    "velocity":
    {
        "type" : "quasineutral" // $U_{\parallel,s} = -\sum_{i,i\neq s}z_i N_i U_{\parallel,i}/(z_sN_s)$
    }
}
\end{minted}
The result of the velocity initialisation is that the initial $A_\parallel=0$ and thus $W_\parallel = U_\parallel$.


\paragraph{Transformation to gyro-center quantities}
Per default we think of the initialisation as determining the physical, particle quantities, however, the
model is derived in terms of gyro-centre quantities.
The physical quantities therefore need to be transformed to gyro-centre quantities. This is in general a
difficult task and may only be possible approximately.
For the densities we allow two different invert modes
\begin{minted}[texcomments]{js}
"init":
{
    "density" :
    {
        "invert" : "none"  // $N_s = n_s$
    },
    "density":
    {
        "invert" : "lwl" // $N_s = n_s - \Delta_\perp \left( \frac{\mu_s p_\perp }{2 z_s^2 B}\right)$
    }
}
\end{minted}
For the remaining quantities we use
\begin{align}
U_\parallel &= u_\parallel\\
P_\perp &= p_\perp = n t_\perp \\
P_\parallel &= p_\parallel = n t_\parallel\\
Q_\perp &= 0 \\
Q_\parallel &= 0
\end{align}


%%%%%%%%%%%%%%%%%%%%%%%%%%%%%%%%%%%%%%%%%%%%%%%%%%%%%%%%%%%%%%%%%%%%%
\subsubsection{Sources} \label{sec:sources}
See the file {\tt init.h} to add your own custom source.

At this stage we can have sources only in the density, perpendicular pressure and
parallel pressure equations. The parallel momentum and heat flux sources are always zero.
\begin{minted}[texcomments]{js}
"source":
{
    "density":
    {
        //...
    },
    "pperp":
    {
        //...
    },
    "ppara":
    {
        //...
    }
}
\end{minted}


\paragraph{Minimum density and temperature}
The first task of the source is to globally ensure a minimum density and temperature. This
is required since through sheath dissipation the density and temperatures can in
principle become arbitrarily close to zero. This would however
be detrimental to the stability of the simulation (divide by zero errors) and in reality
also never happens due to wall-recycling etc.
\begin{align}
    S_{N,s} &= - \omega_{\min} ( N_s - n_{\min,s}) \Theta_{\beta n_{\min}/2} (n_{\min,s} - \beta n_{\min,s} / 2 - N_s)
    \label{eq:minimumne}
    \\
    S_{P} &= -\omega_{\min} N( T - T_{\min,s}) \Theta_{\beta T_{\min}/2}( T_{\min} - \beta T_{\min}/2 - T) + TS_{N,s}
    \label{eq:minimump}
\end{align}
\begin{tcolorbox}[title=Note]
    We do not add Eqs.~\eqref{eq:minimumne} and \eqref{eq:minimump} as a source
    in the file output or in the expression for the velocity source later on.
\end{tcolorbox}
The Heaviside function ensures that this source term only acts when the density
is below the lower limit. $n_{\min}=0$ disables this term.
\begin{minted}[texcomments]{js}
"source":
{
    "minn" : 0.1 // $n_{\min,s}$ in Eq.\eqref{eq:minimumne}
    "mint" : 0.5 // $t_{\min,s}$ in Eq.\eqref{eq:minimump}
    // Note that $\sum_s z_s n_{\min, s} = 0$ must be fulfilled
    //if 0 then the other two parameters are ignored
    "minrate" : 1 // $\omega_{\min}$ in Eq.\eqref{eq:minimumne} (only if $n_{\min}>0$)
    // Must be equal for all species
    "minbeta" : 0.25 // $\alpha_s = \beta n_{\min,s}$ in Eq.\eqref{eq:minimumne} and $\alpha_s = \beta n_{\min,s}$ in Eq.\eqref{eq:minimump} (only if $n_{\min}>0$)
}
\end{minted}
\noindent
In order to disable any additional source terms choose
\begin{minted}[texcomments]{js}
"source":
{
    "density": // or "perp" or "ppara"
    {
        "type" : "zero" // $S_{N} = 0$ or $S_{P,\perp}=0$ or $S_{P,\parallel} = 0$
    }
}
\end{minted}

\paragraph{Profile forcing}
\noindent We can choose the source terms $S_N$, $S_{P,\perp}$ and $S_{P,\parallel}$ to force a profile
\begin{align} \label{eq:electron_source}
    S_{N,s}(R,Z,\varphi, t) &= \omega_N (N_{prof}(R,Z) - N_s(R,Z,\varphi, t))D(R,Z) \\
\end{align}
and analogously for the pressure sources,
where $\omega_N$ is the source strength parameter (Must be equal among species).
The forced source will result in exponential adaption of the core
profile of the form $N \propto N_{prof}+(N_{prof}-N_{0})e^{-\omega_Nt}$.

The source profile is chosen in physical space, i.e. $n_{prof}$, $p_{\perp,prof}$ and $p_{\parallel,prof}$.
\begin{minted}[texcomments]{js}
"source":
{
    "density" : // or "pperp" or "ppara"
    {
        "type" : "fixed_profile", // must be equal for all species and equations
        "rate" : 1e-3 // source strength $\omega_s$
        // Density source rates must be equal among species
        // Pressure source rates can be different
        "profile" :
        {
            // Determine $n_{prof}(R,Z)$
            // see Profile section in `feltor.pdf`
        },
        "damping" :
        {
            // $D(R,Z)$ determines where the source is active
            // see Damping section in `feltor.pdf`
            // The background parameter must be present but is ignored ($n_{bg} = 0$)
        }
    }
}
\end{minted}
To ensure quasineutrality we need to enforce $\sum_s z_s n_{prof,s} = 0$. To achieve
Further, we require all density source rates $\omega_{n,s}$ be equal among species.
In order to achieve this a single species (usually the electrons) can (but does
not have to) be automatically determined by the quasineutrality condition setting
\begin{minted}[texcomments]{js}
"source":
{
    "density" :
    {
        "type" : "fixed_profile-quasineutral"  // $n_{s,prof} = -\sum_{i,i\neq s} z_i n_{i,prof}/ z_s$
        "rate" : 1e-3 // source strength $\omega_s$ must equal other species
        "damping" : // does not have to equal other species
        {
            // $D(R,Z)$ determines where the source is active
            // see Damping section in `feltor.pdf`
            // The background parameter must be present but is ignored ($n_{bg} = 0$)
        }
    },
}
\end{minted}
and transform these to gyrocentre space via
\begin{align}
   N_{prof} &= n_{prof} - \Delta_\perp \left( \frac{\mu_s p_{\perp,prof}}{2z_s^2B^2}\right) - \nc\left( \frac{\mu_s n_{prof} \vn_\perp \phi}{z_sB^2}\right) \\
   P_{\perp,prof} &= p_{\perp,prof} - \nc\left( \frac{\mu_s p_{\perp,prof}\vn_\perp \phi}{z_sB^2}\right) \\
   P_{\parallel,prof} &= p_{\parallel,prof}
\end{align}


\paragraph{Constant influx}
Another option is to provide a constant influx of particles and/or heat (again, in physical space)
\begin{align}
    s_{n,s}(R,Z,\varphi, t) &= \omega_n \left( (n_{prof}(R,Z) + \tilde n(R,Z,\varphi) - n_{bg})D(R,Z) + n_{bg}\right)
\end{align}
where $\omega_n$ is the source strength parameter (equal for all species for density, but can be different for pressure).
The idea is that you can start with zero density and evolve the profile purely
with the source. Destabilization of the profile works if you put a turbulent
bath on top of it.
\begin{tcolorbox}[title=Note]
    % Update to actual output fields
    See the output \mintinline{c}{sne_tt_ifs} in \mintinline{bash}{thermaldiag.cu} for how much mass
the source with the parameters below generates and compare to
\mintinline{c}{jsne_tt_fsa} to see how much mass is lost.
\end{tcolorbox}
\begin{minted}[texcomments]{js}
"source":
{
    "density" : // or "pperp" or "ppara"
    {
        "type" : "influx",
        "rate" : 1e-3, // source strength $\omega_s$
        // Density source rates must be equal among species
        // Pressure source rates can be different
        "ntilde" :
        {
            // Determine $\tilde n(R,Z,\varphi)$
            // see Ntilde section in `feltor.pdf`
            // The turbulent bath works well to destabilize the profile
        },
        "profile" :
        {
            // Determine $n_{prof}(R,Z)$
            // see Profile section in `feltor.pdf`
        },
        "damping" :
        {
            // Determine $D(R,Z)$
            // see Damping section in `feltor.pdf`
            // The boundary determines where the source is active
        }
    }
}
\end{minted}

In order to not generate potential with the constant influx source term the density source
terms need to fulfill $\sum_s z_s s_{n,s}  = 0$. Generally, we achieve this by
requiring all density source rates $\omega_{n,s}$ be equal among species and the source
profile fulfill the quasineutrality condition.
In order to achieve this a single species (usually the electrons) can (but does
not have to) be automatically determined by the quasineutrality condition setting
\begin{minted}[texcomments]{js}
"source":
{
    "density" :
    {
        "type" : "influx-quasineutral"  // $s_{n,s} = -\sum_{i,i\neq s} z_i s_{n,i}/ z_s$
        "rate" : 1e-3 // source strength $\omega_s$ must equal other species
    },
}
\end{minted}

Thereafter we transform the physical sources to gyro-centre space via
\begin{align}
   S_N &= s_n - \Delta_\perp \left( \frac{\mu_s s_{p,\perp}}{2z_s^2B^2}\right) - \nc\left( \frac{\mu_s s_n \vn_\perp \phi}{z_sB^2}\right) \\
   S_{P,\perp} &= s_{p,\perp} - \nc\left( \frac{\mu_s s_{p,\perp}\vn_\perp \phi}{z_sB^2}\right) \\
   S_{P,\parallel} &= s_{p,\parallel}
\end{align}
for all species.

Note that the additional terms besides $s_{n}$ are total divergences which
means they do not change the volume integrated "total" particle number created
by the source (analogous for pressure).  Note that $s_{n}$ needs to be smooth so that $\np^2 s_{n}$
is well defined.

\paragraph{Velocity source}
The velocity source is chosen such that the density source does not introduce
net momentum
\begin{align}
    S_U = -\frac{U_\parallel}{N} S_N
    \label{eq:velocity_source}
\end{align}
for all species. Note that this requires the density source to be transformed to the staggered grid.

%%%%%%%%%%%%%%%%%%%%%%%%%%%%%%%%%%%%%%%%%%
\subsection{Advection}
No options are currently choosable in this section!

We implemented an upwind-advection scheme for the perpendicular terms.
No other scheme is currently choosable.

For the parallel advection we implemented a velocity-staggered finite-volume
scheme. No slope limiter is used.
%%%%%%%%%%%%%%%%%%%%%%%%%%%%%%%%%%%%%%%%%%%%%%%%%%%%%%%%%
\subsection{Regularization} \label{sec:regularization}
In order to stabilize the advection schemes and/or remove cascading energy
from small scale wave components
we implement an artificial viscosity of order $s$ Eq.~\eqref{eq:perpdiffNU}
for the perpendicular terms and artificial viscosity of order 1 of the parallel terms.
\begin{tcolorbox}[title=Note]
    We apply the viscosity to the parallel velocities $U_\parallel$ and
    not to the parallel canonical velocities $W_\parallel$.
\end{tcolorbox}
\begin{minted}[texcomments]{js}
"regularization",
{
    "order" : 2,  // order of the perpendicular artifical diffusion
    "direction" : "centered",  // The direction of the Laplacian
    "nu_perp" : [1e-4, 1e-4, 1e-4, 1e-4, 1e-4, 1e-4],
    // The strength of the perpendicular diffusion for density, tperp, tpara, ...
    "nu_parallel" : [500, 500, 500, 500, 500, 500],
    // The strength of the parallel diffusion for density, tperp, tpara, ...
}
\end{minted}
Simply set the diffusion to 0 if you do not want any regularization.
\begin{tcolorbox}[title=Note]
    Having to make the timestep smaller due to viscosity
    is a clear warning that \mintinline{c}{nu_perp} is too large for the
    chosen resolution.
\end{tcolorbox}
%%%%%%%%%%%%%%%%%%%%%%%%%%%%%%%
\subsection{Elliptic solvers}
The section `Elliptic solvers` in `feltor.pdf` applies verbatim.

%%%%%%%%%%%%%%%%%%%%%%%%%%%%%%%%%%%%%%%%%
\subsection{FCI: Parallel derivatives}
The section `FCI: Parallel derivatives` in `feltor.pdf` applies verbatim.

%%%%%%%%%%%%%%%%%%%%%%%%%%%%%%%%%%%%%%%%%%%%%%%%%%%%%%%%%
\subsection{Parameters in equations} \label{sec:physical}
\begin{minted}[texcomments]{js}
"physical":
{
    "beta" : 5e-5, // Plasma beta $5\cdot 10^{-6}$ (TJK), $4\cdot 10^{-3}$ (Compass),
    // If $0$, then the model is electrostatic
    "collisionality" : 1e-4, // Rerefence collisionality $\nu_{ref}$
    "mue" : 2.72121e-4, // $\mu_e =m_e/m_i$. [IGNORED]  This parameter is
    // ignored by all codes, its sole purpose is to make the set of numerical
    // parameters uniquely invertible for analysis purposes (see discussion below)
    "epsilon_D" : 1e-5, // [IGNORED] Eq.\eqref{eq:epsilon_D} This parameter is
    // ignored by all codes, its sole purpose is to make the set of numerical
    // parameters uniquely invertible for analysis purposes (see discussion below)
    "qlandau" : 3, // Parameter for Landau damping in heat flux equations
}
\end{minted}
These parameters, together with the $R_0$ parameter in the magnetic field,
determine the physical parameter regime we are in.
In total, our model has \textbf{4} independent dimensionless parameters
\begin{align*}
    \mu_e = \frac{m_e}{m_i}\ \beta =
    \frac{n_0T_{e}}{B_0^2/\mu_0}\ \nu_{ref} = \frac{\sqrt{2m_e} e^3 \ln
    \Lambda} {12\pi^{3/2} \epsilon_0^2} \frac{n_0} {B_0T_e^{3/2}} \ R_0
    = \frac{\hat R_0 e B_0}{\sqrt{T_{e}m_i}}
\end{align*}
These are computed from \textbf{5} physical parameters
\begin{align*}
    m_i\ T_e\ n_0\ B_0\ \hat R_0
\end{align*}
 A given set of numerical parameters thus represents a one-dimensional
 submanifold of physical parameters.
 This is because we neglect the Debye length relative to the ion gyro-radius
 in the model
 \begin{align}\label{eq:epsilon_D}
     \epsilon_D = (\lambda_D/\rho_s)^2 = \epsilon_0 T_e /n_0 /e^2 / T_e/m_i e^2 B_0^2 = \epsilon_0 B_0^2/n_0 /m_i \rightarrow 0
 \end{align}
If we further neglect $\beta=0$ this becomes a two-dimensional submanifold
of physical parameters.
\begin{tcolorbox}[title=Note]
    In order to invert the system (for given numerical parameters find the
    corresponding physical parameters) we have to (i) give $\epsilon_D$ in the
    input file or (ii) fix one (or more) physical
    parameters to get a solvable system.  In practice this may be the scale of
    the machine $\hat R_0$ (and after that the magnetic field strength $B_0$).

    Note that $\mu_e$ in the above list is ignored because the species mass ratio is given
    in the species list.
\end{tcolorbox}
%%%%%%%%%%%%%%%%%%%%%%%%%%%%%%%%%%%%%%%%%%%%%%%%%%%%
\subsection{Timestepper}
The section `Timestepper` in `feltor.pdf` applies verbatim.

\subsection{Steering the file output} \label{sec:output}
The section `Steering the file output` in `feltor.pdf` applies verbatim except
that "GLFW" output is not possible.

%%%%%%%%%%%%%%%%%%%%%%%%%%%%%%%%%%%%%

\subsection{[Optional] Output equations flags}\label{sec:output_flags}
The optional variable "equations" contains a list of strings that
determine which variables will be written into the output file.
Each string
corresponds to a different equation that can be studied with FELTOR at the
moment. These equations are described in section \ref{sec:output}, and their
definitions are in the equations titles.
\begin{minted}[texcomments]{js}
"output":
{
    "equations" :
    [
        "Basic",
        "Mass-conserv",
        "Energy-theorem",
        "Toroidal-momentum",
        "Parallel-momentum",
        "Zonal-Flow-Energy",
        "COCE"
    ]
}
\end{minted}
\begin{tcolorbox}[title=Note]
    If the "equations" field does not exist, the default selection is ["Basic", "Mass-conserv", "Energy-theorem","Toroidal-momemtum","Parallel-momentum","Zonal-Flow-Energy"]
\end{tcolorbox}
%%%%%%%%%%%%%%%%%%%%%%%%%%%%%%%%%%%%%%%%%%%%%%%

\subsection{[Optional] Probes (synthetic diagnostics)}
The section `[Optional] Probes (synthetic diagnostics)` in `feltor.pdf` applies verbatim
except that the first sentence applies to $N$, $U_\parallel$, $T_\perp$, $T_\parallel$, $Q_\perp$ and $Q_\parallel$ for all species s.

%%%%%%%%%%%%%%%%%%%%%%%%%%%%%%%%%%%%%%%%%%%%%%%
\subsection{[Optional] Computational flags}
We can simplify the full simulation in a few ways:
\begin{minted}[texcomments]{js}
"flags": [
    "symmetric", // assume toroidal symmetry, in this mode the
    // equations become 2 dimensional and the third dimension in the init
    // condition is ignored, the Nz parameter in the grid is used to initialize
    // the FCI method and then overwritten to $N_z\equiv 1$.
    "calibrate", // Simulation stops after initialization,
    // no right hand side calls
    // not possible for glfw output
    "isothermal", // $T_\perp = T_\parallel = 1$, $Q_\perp = Q_\parallel = 0$, $s_{p_\perp} = s_{p_\parallel} = s_n$
]
\end{minted}
You can combine the flags in any way you want: you can have 1 or more
flags present or leave the field empty.
\begin{tcolorbox}[title=Note]
    Use ["symmetric", "calibrate"] together (to avoid FCI initialization
    which takes a long time) and quickly fine-tune
    magnetic field, wall, sheath, init, source and grid parameters
\end{tcolorbox}
%%%%%%%%%%%%%%%%%%%%%%%%%%%%%%%%%%%%%%%%%%%%%%%%%%
\subsection{Restarting} \label{sec:restart}
The section `Restarting` in `feltor.pdf` applies verbatim except that the code is called `thermal` instead of `feltor` and the computational fields include
$N$, $U_\parallel$, $T_\perp$, $T_\parallel$, $Q_\perp$ and $Q_\parallel$ for all species s.
Also, all species names, mass charge, kappa and pi values in the input file must be identical to the ones in the restart file.

%%%%%%%%%%%%%%%%%%%%%%%%%%%%%%%%%%%%%%%%%%%%%
\section{Output} \label{sec:output_file}
Output file format: \href{https://www.unidata.ucar.edu/software/netcdf/docs/}{netcdf-4/hdf5};

A \textit{coordinate variable (Coord. Var.)} is a Dataset with the same name as a dimension.
We follow
\href{http://cfconventions.org/Data/cf-conventions/cf-conventions-1.7/cf-conventions.html}{CF Conventions CF-1.7}
and write according attributes into the file.

\begin{tcolorbox}[title=Note]
    The command \mintinline{bash}{ncdump -h output.nc} gives a full list of what a file contains.
\end{tcolorbox}
Here, we list the content without attributes
since the internal netcdf information does not display equations.
%
%Name | Type | Dimensionality | Description
%---|---|---|---|
\begin{longtable}{lll>{\RaggedRight}p{7cm}}
\toprule
\rowcolor{gray!50}\textbf{Name} &  \textbf{Type} & \textbf{Dimension} & \textbf{Description}  \\ \midrule
inputfile  &     text attribute & - & verbose input file as a string (valid JSON, C-style comments are allowed but discarded) \\
x                & Coord. Var. & 1 (x) & $R$-coordinate (computational space, compressed size: $nN_x/c_x$)\\
y                & Coord. Var. & 1 (y) & $Z$-coordinate (computational space, compressed size: $nN_y/c_y$)\\
z                & Coord. Var. & 1 (z) & $\varphi$-coordinate (computational space, size: $N_z$) \\
time             & Coord. Var. & 1 (time)& time at which fields are written (variable size: maxout$+1$, dimension size: unlimited) \\
xc           & Dataset & 3 (z,y,x) & Cartesian x-coordinate $x=R\sin(\varphi)$ \\
yc           & Dataset & 3 (z,y,x) & Cartesian y-coordinate $y=R\cos(\varphi)$\\
zc           & Dataset & 3 (z,y,x) & Cartesian z-coordinate $z=Z$ \\
Psip             & Dataset & 3 (z,y,x) & Flux function $\psi_p(R,Z)$ \\
vol3d            & Dataset & 3 (z,y,x) & Volume form in 3d including dG weights (can be used to integrate 3d quantities)  \\
Nprof            & Dataset & 3 (z,y,x) & Density profile $n_\text{prof}$ used in the forcing source \\
Source           & Dataset & 3 (z,y,x) & Source profile $S_{prof}$\\
BR               & Dataset & 3 (z,y,x) & Contravariant magnetic field component $B^R$ \\
BZ               & Dataset & 3 (z,y,x) & Contravariant magnetic field component $B^Z$ \\
BP               & Dataset & 3 (z,y,x) & Contravariant magnetic field component $B^\varphi$ \\
electrons        & Dataset & 4 (time, z, y, x) & electron density $n_e$ \\
ions             & Dataset & 4 (time, z, y, x) & ion density $N_i$ \\
Ue               & Dataset & 4 (time, z, y, x) & electron velocity $u_{\parallel,e}$ \\
Ui               & Dataset & 4 (time, z, y, x) & ion velocity $U_{\parallel,i}$ \\
potential        & Dataset & 4 (time, z, y, x) & electric potential $\phi$ \\
aparallel        & Dataset & 4 (time, z, y, x) & parallel vector potential $A_\parallel$ \\
X\_2d            & Dataset & 3 (time,y,x) & Selected plane $X(\varphi=0)$ \\
X\_ta2d          & Dataset & 3 (time,y,x) & Toroidal average $\PA{ X }$
Eq.~\eqref{eq:phi_average} \\
Y\_tt\_2d        & Dataset & 3 (time,y,x) & Time integrated (between two outputs, Simpson's rule) selected plane
$\int_{t_0}^{t_1}\d t Y(\varphi=0) $
where $t_1 - t_0 = \Delta T_{\text{out}}$, $t_1$ is the current time and $t_0$ is the time of the last output. \color{red} The first output at time $t_1=0$ is the same as X\_2d !!\color{black} \\
Y\_tt\_ta2d      & Dataset & 3 (time,y,x) & Time integrated (between two outputs, Simpson's rule) toroidal average (Eq.~\eqref{eq:phi_average})
$\int_{t_0}^{t_1}\d t \PA{ Y }$
where $t_1 - t_0 = \Delta T_{\text{out}}$, $t_1$ is the current time and $t_0$ is the time of the last output. \color{red} The first output at time $t_1 = 0$ is the same as X\_2d !!\color{black} \\
\bottomrule
\end{longtable}
where
X and Y\_tt represent the quantities described in the tables in this section starting with the quantities
\begin{longtable}{llll}
\toprule
\rowcolor{gray!50}\textbf{Name} &  \textbf{Equation} & \textbf{Name} &  \textbf{Equation}\\
\midrule
    electrons &$n_e$ &
    ions &$N_i$ \\
    Ue &$u_{\parallel,e}$ &
    Ui &$U_{\parallel,i}$ \\
    potential &$\phi$ &
    psi &$\psi$ \\
    aparallel &$A_\parallel$ & \\
    gammaNi & $\Gamma_{1,i} N_i$  &
    gammaPhi & $\Gamma_{1,i} \phi$  \\
    vorticity &$-\Delta_\perp\phi$ &
    vorticity\_i & $-\Delta_\perp \psi$ \\
    laplace\_ne &$\Delta_\perp n_e$ &
    laplace\_ni & $\Delta_\perp N_i$ \\
     & &
    apar\_vorticity &$-\Delta_\perp A_\parallel$ \\
    dssue & $\npar^2 u_{\parallel,e}$&
    %dppue & $\partial_\varphi^2 u_{\parallel,e}$\\
    %dpue2 & $(\partial_\varphi u_{\parallel,e})^2$&
    lperpinv &$L_\perp^{-1} := |\vec\np n_e|/n_e$ \\
    perpaligned &$(\vec\np n_e)^2/n_e$ &
    lparallelinv &$L_\parallel^{-1} := |\npar n_e|/n_e$ \\
    aligned &$ (\npar n_e)^2/n_e$ &
    ne2 & $n_e^2$ \\
    phi2 & $\phi^2$ &
    nephi & $n_e\phi$ \\
\bottomrule
\end{longtable}
The computation time spent on diagnostics is usually negligible unless {\tt itstp}
is unreasonably high. Also remember that the X and Y fields are all
two-dimensional, which takes up much less disk-space than three-dimensional
fields. This outputs are saved if the flag "output":"equations":"Basic" is true.

%\subsection{Conservation laws} \label{sec:conservation}
%In the output file we have
\subsection{Mass conservation, "Mass-conserv"}
The density equations directly yield the particle conservation
\begin{align} \label{eq:mass_theorem}
  \frac{\partial}{\partial t} N
  + \nc\vec{ j_{N}}
  =  \Lambda_{N}+S_{N}
\end{align}
The terms of the particle conservation thus read
\begin{align}
  N= & N,\\
  \vec j_{N} =& N\left(
  \vec u_\psi + \vec u_C + \vec u_{K} +U_\parallel\left(\bhat+{\vec b}_\perp\right)  \right)
\label{eq:particle_flux}\\
  %\nonumber\\
  %=& N \left(\frac{\bhat\times \vn\phi}{B}
  %+ \tau_e \frac{\bhat\times\vn n_e}{n_eB}
  %+ \mu_e u_{\parallel,e}^2\vec K_{\vn\times\bhat}
  %+ u_{\parallel,e}(\bhat + {\vec b}_\perp) \right), \\
  \Lambda_{N} =& \Lambda_N
\\
  S_{N} =&  S_{N}
\end{align}
Notice that
\begin{align}
\tau N \vec K = \tau N\vn\times\frac{\bhat}{B} = \tau \vn\times N\frac{\bhat}{B} + \tau \frac{\bhat\times\vn N}{B}
\label{}
\end{align}
such that we can define the diamagnetic flux in the particle flux since
the rotation vanishes under the divergence.

We here also derive the particle flux \eqref{eq:particle_flux} through a flux surface
\begin{align} \label{eq:radial_particle_flux}
 \vec j_{N}\cn v %=& N\left( \vec u_E + \vec u_C + \vec u_{\vn
 %B} + U_\parallel \left(\bhat + {\vec b}_\perp\right)\right) \cn \psi_p \nonumber\\
 =&
  \frac{\d v}{\d \psi_p} N\left[\frac{1}{B}[\psi, \psi_p]_\perp + \left(\tau + \mu U_\parallel^2\right)
   \mathcal K_{\vn\times\bhat}(\psi_p) + \tau  \mathcal K_{\vn B}(\psi_p) \right] \nonumber\\
 &+ NU_\parallel\frac{\d v}{\d \psi_p}\left [\left( A_\parallel \mathcal
 K_{\vn\times\bhat}(\psi_p) + \frac{1}{B}[\psi_p, A_\parallel]_\perp\right) \right]
\end{align}
The relevant terms in the output file are
\begin{longtable}{llll}
\toprule
\rowcolor{gray!50}\textbf{Name} &  \textbf{Equation} & \textbf{Name} &  \textbf{Equation}\\
\midrule
    electrons & $n_e$ &
    jscurvne\_tt &$ n_e  \vec u_K \cn \psi_p$ \\
    jscurvkappane\_tt &$ n_e  \vec u_C \cn \psi_p$ &
    & \\
    jsneA\_tt &$ n_e u_{\parallel,e} \vec{ b}_\perp  \cn \psi_p$ &
    jsneE\_tt & $ n_e \vec u_E\cn\psi_p$ \\
    lneperp\_tt &$ \Lambda_{\perp,n_e} = -\nu_{n,\perp} (-\Delta_\perp)^s n_e$&
    lneparallel\_tt &$ \Lambda_{\parallel,n_e} = \nu_{N\parallel} \Delta_\parallel n_e$ \\
    sne\_tt & $S_{n_e}$ &
    jsdiae\_tt & $\tau_e \bhat \times \vn n_e \cn \psi_p /B$\\
    divjnepar\_tt & $\nc (\bhat n_e u_{e,\parallel}$) &
    divneE\_tt & $\nc ( n_e \vec u_E)$ \\
    divcurvne\_tt & $\nc( \tau_e \vec K n_e)$ &
    divcurvkappane\_tt & $\nc( \mu_e n_e u_{\parallel,e}^2 \KK)$ \\
    divneA\_tt & $\nc (\bhat n_e u_{e,\parallel} {\vec b}_\perp)$) &
     &  \\
    ions & $N_i$ &
    jscurvni\_tt &$ N_i \vec u_K \cn \psi_p$ \\
    jscurvkappani\_tt &$ N_i \vec u_C \cn \psi_p$ &
    & \\
    jsniA\_tt &$ N_i U_{\parallel,i} \vec{ b}_\perp  \cn \psi_p$ &
    jsniE\_tt & $ N_i \vec u^i_E\cn\psi_p$ \\
    lniperp\_tt &$ \Lambda_{\perp,N_i} = -\nu_{n,\perp} (-\Delta_\perp)^s N_i$&
    lniparallel\_tt &$ \Lambda_{\parallel,N_i} = \nu_{N,\parallel} \Delta_\parallel N_i$ \\
    sni\_tt & $S_{N_i}$ &
    jsdiai\_tt & $\tau_i \bhat \times \vn N_i \cn \psi_p /B$\\
    divjnipar\_tt & $\nc (\bhat N_i U_{i,\parallel}$) &
    divniE\_tt & $\nc ( N_i \vec u_{E,i})$ \\
    divcurvni\_tt & $\nc( \tau_i \vec K N_i)$ &
    divcurvkappani\_tt & $\nc( \mu_i N_i U_{\parallel,i}^2 \KK)$ \\
    divniA\_tt & $\nc (\bhat N_i U_{i,\parallel} {\vec b}_\perp)$) &
     &  \\
\bottomrule
\end{longtable}

Note that the parallel divergences vanish exactly under a flux-surface average. This can serve as a numerical test of our implementation.
\subsection{Energy theorem, "Energy-theorem"}
The terms of the energy theorem are
\begin{align} \label{eq:energy_theorem}
\partial_t \mathcal E +
\nc \vec j_{\mathcal E}
= \Lambda_{\mathcal E}
+  S_{\mathcal E}
+  R_{\mathcal E}
\end{align}
with ( $z_e=-1$ and $z_i=+1$) and $\vec u_E := {\bhat\times \vn\phi}/{B}$
\begin{align} \label{eq:energy_conservation}
  \mathcal{E}= & z_e\tau_e n_e \ln{(n_e)} +z_i\tau_i N_i\ln{(N_i)}
  +\frac{1}{2\beta}\left(\np A_\parallel\right)^2
   +  \frac{1}{2} z_i \mu_i N_i u_E^2  \nonumber\\
   & +\frac{1}{2} z_e\mu_e  n_e u_{\parallel,e}^2
  +\frac{1}{2} z_i\mu_i  N_i U_{\parallel,i}^2,\\
  \vec j_{\mathcal E} =& \sum_s z\left[
  \left(\tau \ln N + \frac{1}{2}\mu U_\parallel^2 + \psi \right)N\left(
  \vec u_E + \vec u_C + \vec u_{K} +U_\parallel\left(\bhat+{\vec b}_\perp\right)  \right) \right]
  \nonumber\\
  &+ \sum_z z\left[\mu \tau NU_\parallel^2\vec K_{\vn\times\bhat} + \tau NU_\parallel \left(\bhat + {\vec b}_\perp\right)\right], \\
  \Lambda_{\mathcal E} =&  \sum_s z\left[\left( \tau\left( 1+\ln{N}\right) + \psi + \frac{1}{2} \mu U_\parallel^2 \right)
  \Lambda_N  +  \mu NU_\parallel\Lambda_U  \right]
\nonumber \\
  S_{\mathcal E} =&  \sum_s  z\left[ \left(\tau\left( 1+\ln{N}\right) +\psi - \frac{1}{2} \mu U_\parallel^2 \right)S_{N}\right]
\nonumber \\
R_{\mathcal E} =&  -\eta_\parallel  n_e(U_{\parallel,i}-u_{\parallel,e})(N_iU_{\parallel,i} - n_eu_{\parallel,e}).
\end{align}
where in the energy flux $\vec j_{\mathcal E}$
we neglect terms  containing time derivatives
of the eletric and magnetic potentials and we sum over all species.
The energy density $\mathcal E$ consists of the Helmholtz free energy density for electrons and ions,
the \(\vec{E} \times \vec{B}\) energy density, the parallel energy densities for electrons and ions and the perturbed magnetic field energy density.
In \(\Lambda\) we insert the dissipative terms of Section~\ref{sec:dissres}. \\
Replace $\Delta_\perp$ with $-\Delta_\perp^2$ when hyperviscous diffusion is chosen
for the diffusion terms in the above equations.

We have the energy flux through a flux surface
\begin{align}
 \vec j_{\mathcal E}\cn v =&%\frac{\d v}{\d \psi_p} \vec j_{\mathcal E}\cn \psi_p  =
\frac{\d v}{\d \psi_p}\sum_s z\left (\tau\ln N + \frac{1}{2}\mu U_\parallel^2 + \psi\right) \vec j_N\cn\psi_p
+ z \mu\tau NU_\parallel^2 \mathcal K_{\vn\times\bhat}(\psi_p) \nonumber\\
&+ z \tau NU_\parallel
 \left( A_\parallel \mathcal
 K_{\vn\times\bhat}(\psi_p) + \frac{1}{B}[\psi_p, A_\parallel]_\perp\right)
\label{eq:energy_flux}
\end{align}
The relevant terms in the output file are
\begin{longtable}{ll}
\toprule
\rowcolor{gray!50}\textbf{Name} &  \textbf{Equation}\\
\midrule
    nelnne &$ z_e\tau_e n_e \ln n_e$ \\
    nilnni &$ z_i\tau_i N_i \ln N_i$ \\
    aperp2 &$ (\np A_\parallel)^2/2/\beta$ \\
    ue2   &$z_i\mu_i N_i u_E^2 /2$ \\
    neue2 &$ z_e\mu_e n_e u_{\parallel,e}^2/2$ \\
    niui2 &$ z_i\mu_i N_i U_{\parallel,i}^2/2$ \\
    see\_tt & $z_e(\tau_e (1+\ln n_e) + \phi + \frac{1}{2}\mu_e u_{\parallel,e}^2) S_{n_e} $ \\
    sei\_tt & $z_i(\tau_i (1+\ln N_i) + \psi + \frac{1}{2}\mu_i U_{\parallel,i}^2) S_{N_i} $ \\
    resistivity\_tt &-$\eta_\parallel n_e (U_{\parallel,i}-u_{\parallel,e})(N_iU_{\parallel,i}-n_eu_{\parallel,e})$ \\
    jsee\_tt &$z_e(\tau_e \ln n_e + \mu_e u_{\parallel,e}^2/2 + \phi)n_e(\vec u_E + \vec u_C + \vec u_K)\cn \psi_p
        + z_e \tau_e n_e u_{\parallel,e}^2 \vec K_{\vn\times\bhat}\cn \psi_p$ \\
    jsei\_tt &$z_i(\tau_i \ln N_i + \mu_i U_{\parallel,i}^2/2 + \psi_i)N_i(\vec u_E^i + \vec u_C + \vec u_K)\cn \psi_p
        + z_i \tau_i N_i U_{\parallel,i}^2 \vec K_{\vn\times\bhat}\cn \psi_p$ \\
    jseea\_tt &$z_e(\tau_e \ln n_e + \mu_e u_{\parallel,e}^2 + \phi)n_e \vec { b}_\perp\cn \psi_p
        + z_e \tau_e n_e u_{\parallel,e} \vec{ b}_\perp \cn \psi_p $ \\
    jseia\_tt &$z_i(\tau_i \ln N_i + \mu_i U_{\parallel,i}^2 + \psi_i)N_i \vec { b}_\perp\cn \psi_p
        + z_i \tau_i N_i U_{\parallel,i} \vec{ b}_\perp \cn \psi_p $ \\
    divee\_tt &$\nc \vec j_{\mathcal E}^e$ \\
    divei\_tt &$\nc \vec j_{\mathcal E}^i$ \\
    diveea\_tt &$\nc \vec j_{\mathcal E,A}^e$\\
    diveia\_tt &$\nc \vec j_{\mathcal E,A}^i$\\
    leeperp\_tt &$z_e(\tau_e(1+\ln n_e) + \phi + \mu_eu_{\parallel,e}^2/2) \Lambda_{n_e,\perp} + z_e\mu_e n_e u_{\parallel,e} \Lambda_{u_e,\perp}$ \\
    leiperp\_tt &$z_i(\tau_i(1+\ln N_i) + \psi_i + \mu_iU_{\parallel,i}^2/2) \Lambda_{N_i,\perp} + z_i\mu_i N_i U_{\parallel,i} \Lambda_{U_i,\perp}$ \\
    leeparallel\_tt & $z_e(\tau_e(1+\ln n_e) + \phi + \mu_eu_{\parallel,e}^2/2) \Lambda_{n_e,\parallel} +
    z_e\mu_e u_{\parallel,e} \Lambda_{u_e, \parallel}$ \\
    leiparallel\_tt & $z_i(\tau_i(1+\ln N_i) + \psi_i + \mu_iU_{\parallel,i}^2/2)\Lambda_{N_i,_\parallel} +
    z_i\mu_i U_{\parallel,i} \Lambda_{U_i,\parallel}$ \\
    %leeparallel\_tt & $ z_e\mu_e u_{\parallel,e} \nu_{\parallel,e} \Delta_\parallel u_{\parallel,e}$ \\
    %leiparallel\_tt & $ z_i\mu_i U_{\parallel,i} \nu_{\parallel,i} \Delta_\parallel U_{\parallel,i}$ \\
    divjeepar\_tt &$ z_e\nc \left( (\tau_e(1+\ln n_e) + \phi + \mu_eu_{\parallel,e}^2/2) n_eu_{\parallel,e}\bhat\right)$ \\
    divjeipar\_tt &$ z_i\nc \left( (\tau_i(1+\ln N_i) + \psi + \mu_iU_{\parallel,i}^2/2) N_iU_{\parallel,i}\bhat\right)$ \\
\bottomrule
\end{longtable}

\subsection{Toroidal ExB angular momentum equation, "Toroidal-momentum"} \label{sec:vorticity_eq}
We integrate the polarisation equation over volume, multiply by $\d \psi_p/\d v$ and derive by time. In the drift-ordering up to order $\mathcal O(\delta^3)$ we get
\begin{align}
    &\partial_t \RA{\Omega} + \frac{\partial}{\partial v}\frac{\d v}{\d\psi_p}\RA{\vec j_\Omega\cn\psi_p} = -\RA{F_{L,\varphi}} + \RA{\mathcal S_\Omega} \label{eq:vorticity_average} \\
\Omega &:= \mu_i N_i \left(\frac{\vn\psi_p\cn\phi}{B^2} + \tau_i \vn\ln N_i\cn\psi_p\right) \equiv \mu_i N_i(u_{E,\varphi} + u_{D,\varphi}) \\
\vec j_{\Omega} &:= \Omega \vec u_E
    - \left(\frac{1}{\beta} \vn\psi_p \cn A_\parallel +\frac{1}{2}\tau_i \vn\psi_p\cn  (N_iU_{\parallel,i})\right)\frac{\bhat\times\vn A_\parallel}{B} \\
    F_{L,\varphi} &:=  -(z_e \tau_e n_e + z_i\tau_i N_i)\mathcal K(\psi_p) - (z_e\mu_e n_eu_{\parallel,e}^2 + z_i\mu_i N_iU_{\parallel,i}^2)\mathcal K_{\vn\times\bhat}(\psi_p) \\
    \mathcal S_\Omega &:= \mu_i S_{n_e} \frac{\vn\psi_p\cn \phi}{B^2} + \mu_i\tau_i\vn\psi_p\cn S_{n_e} \label{eq:em_source}
\end{align}
Equation~\eqref{eq:vorticity_average} can be rewritten by inserting the continuity equation to yield an equation only for the \ExB angular momentum. Again up to order $\mathcal O(\delta^3)$ in the drift ordering we obtain
(the diffusive term is for testing purposes)
\begin{align}
&\partial_t \RA{\Omega_E} + \frac{\partial}{\partial v} \frac{\d v}{\d \psi_p}\RA{ \vec j_{\Omega_E}\cn\psi_p} = -\RA{F_{L,\varphi}}+ \RA{\mathcal S_{\Omega_E}} + \RA{\Lambda_{\Omega_E}} \label{eq:exb_average} \\
\Omega_E &:= \mu_i N_i \frac{\vn\psi_p\cn\phi}{B^2} \equiv \mu_i N_i u_{E,\varphi} \\
\vec j_{\Omega_E} &:= \Omega_E (\vec u_E + \vec u_D)
    - \vn A_\parallel\cn\psi_p \left(\frac{1}{\beta} \frac{\bhat\times\vn A_\parallel}{B} +\frac{1}{2} \bhat \times \vn \mu_i \tau_i N_iU_{\parallel,i}\right) \\
    \mathcal S_{\Omega_E} &:= \mu_i S_{n_e} \frac{\vn\psi_p\cn\phi}{B^2} \quad
    \Lambda_{\Omega_E} := \mu_i \Lambda_{n_e}\frac{\vn\psi_p\cn\phi}{B^2}
\end{align}
where here we also monitor the source and diffusion terms.
In the output file we have
\begin{longtable}{llll}
\toprule
\rowcolor{gray!50}\textbf{Name} &  \textbf{Equation}&
\textbf{Name} &  \textbf{Equation}\\
\midrule
    oexbi &$\mu_i N_i \frac{\vn\psi_p\cn\phi}{B^2}$ &
    oexbe &$\mu_i n_e \frac{\vn\psi_p\cn\phi}{B^2}$ \\
    odiai &$\mu_i \tau_i\vn\psi_p\cn N_i$ &
    odiae &$\mu_i \tau_i\vn\psi_p\cn n_e$ \\
    divoexbi\_tt &$\mu_i \nc( N_i \frac{\vn\psi_p\cn\phi}{B^2} \vec u_E)$ &
    divoexbe\_tt &$\mu_i \nc( n_e \frac{\vn\psi_p\cn\phi}{B^2} \vec u_E)$ \\
    divoexbiUD\_tt &$\mu_i\tau_i \nc( \frac{\vn\psi_p\cn\phi}{B^2} \frac{\bhat\times\vn N_i}{B})$ &
    divoexbeUD\_tt &$\mu_i\tau_i \nc( \frac{\vn\psi_p\cn\phi}{B^2} \frac{\bhat\times\vn n_e}{B})$ \\
    divodiaiUE\_tt &$\mu_i \tau_i\nc(\vn\psi_p\cn N_i \vec u_E)$ &
    divodiaeUE\_tt &$\mu_i \tau_i\nc(\vn\psi_p\cn n_e \vec u_E)$ \\
    divoApar\_tt &$ -\nc ( \vn\psi_p\cn A_\parallel \frac{\bhat\times\vn A_\parallel}{B\beta})$ &
    & \\
    jsoexbi\_tt &$\mu_i N_i \frac{\vn\psi_p\cn\phi}{B^2} \frac{\bhat\times\vn\phi\cn \psi_p}{B}$ &
    jsoexbe\_tt &$\mu_i n_e \frac{\vn\psi_p\cn\phi}{B^2} \frac{\bhat\times\vn\phi\cn \psi_p}{B}$ \\
    jsodiaiUE\_tt &$\mu_i \tau_i\vn\psi_p\cn N_i \frac{\bhat\times\vn\phi\cn \psi_p}{B}$ &
    jsodiaeUE\_tt &$\mu_i \tau_i\vn\psi_p\cn n_e \frac{\bhat\times\vn\phi\cn \psi_p}{B}$ \\
    jsoexbiUD\_tt &$\mu_i\tau_i \frac{\vn\psi_p\cn\phi}{B^2} \frac{\bhat\times\vn N_i\cn \psi_p}{B}$ &
    jsoexbeUD\_tt &$\mu_i\tau_i \frac{\vn\psi_p\cn\phi}{B^2} \frac{\bhat\times\vn n_e\cn \psi_p}{B}$ \\
    jsoApar\_tt &$ -\vn\psi_p\cn A_\parallel \frac{\bhat\times\vn A_\parallel\cn \psi_p}{B\beta}$ &
    jsodiaApar\_tt & $ -\frac{1}{2}\tau_i \vn\psi_p\cn  (N_iU_{\parallel,i})\frac{\bhat\times\vn A_\parallel}{B}\cn\psi_p$ \\
    jsoexbApar\_tt & $ -\frac{1}{2}\tau_i \bhat\times\vn  (N_iU_{\parallel,i})\cn\psi_p \vn A_\parallel\cn\psi_p$ &
    & \\
    socurve\_tt &$z_e\tau_e n_e \mathcal K(\psi_p)$ &
    socurvkappae\_tt &$z_e\mu_e n_eu_{\parallel,e}^2 \mathcal K_{\vn\times\bhat}(\psi_p)$ \\
    socurvi\_tt &$z_i\tau_i N_i \mathcal K(\psi_p)$ &
    socurvkappai\_tt &$z_i\mu_i N_iU_{\parallel,i}^2 \mathcal K_{\vn\times\bhat}(\psi_p)$ \\
    sosne\_tt & $\mu_i S_{n_e} \vn\psi_p\cn\phi/B^2$ &
    sospi\_tt & $\mu_i \tau_i \vn\psi_p \cn S_{n_e}$\\
    loexbe\_tt & $ \mu_i \Lambda_{n_e} \vn\psi_p\cn\phi/B^2$ & \\
\bottomrule
\end{longtable}

%%%%%%%%%%%%%%%%%%%%%%%%%%%%%%%%%%%%%%%
\subsection{Conservation of Currents (COCE), "COCE"}\label{sec:COCE}
If we want to study the electric field locally, without the flux surface average implemented in the equation at section \ref{sec:vorticity_eq}, the COCE diagnostic is necessary. The equation implemented is 
\begin{equation}
    \partial_t\Omega-\nc\nc(\vec{\omega u_E})=\nc(j_\parallel\bhat+\vec{j_{curv}}+\vec{j_{b_\perp}})+\Omega_S,
    \label{eq:COCE_diagnostic}
\end{equation}   
with the following definitions:
\begin{equation}
\Omega=\Omega_E+\Omega_D=-\nc\vec{\omega}=-\nc(\vec{\omega_D}+\vec{\omega_E})=\nc\left[\frac{m_i}{B^2}\left(\frac{\nabla_\perp p_i}{q_i}+n_i\nabla_\perp\phi\right)\right],
\end{equation}
\begin{equation}
    \Omega_S=\nc\left[\frac{m_i}{B^2}\left(\frac{\nabla_\perp S_{p_i}}{q_i}+S_{n_i}\nabla_\perp\phi\right)\right],
\end{equation}
\begin{equation}
     \vec{j_{curv}}\approx\frac{\bhat\times\nabla(p_e+p_i)}{B}+(j_\parallel A_{1,\parallel}+\sum_s m_s n_s u_{\parallel,s}^2) \frac{\nabla\times\bhat}{B}=\vec{j_{d}}+\vec{j_A}+\vec{j_\kappa},
\end{equation}
\begin{equation}
      \vec{j_{b_\perp}}=(j_\parallel+j_{mag})\vec{b_\perp}-\nc(\vec{M^{em}}\vec{b_\perp})=j_\parallel\vec{b_\perp}+\vec{j_{M^{em}}}.
\end{equation}

The implemented terms are defined as follows:
\begin{longtable}{ll}
\toprule
\rowcolor{gray!50}\textbf{Name} &  \textbf{Equation}\\
\midrule
    v\_Omega\_E\_gf &$ \mu_i\nc(\frac{N_i\vec{\nabla_\perp}\phi}{B^2})$ \\
    v\_Omega\_D\_gf &$ \mu_i\tau_i \Delta_\perp N_i$ REMOVED! use laplace\_ne and laplace\_ni \\
    v\_adv\_E\_main\_gf\_tt &$ \mu_i\nc(\nc (\frac{N_i\vec{\nabla_\perp}\phi}{B^2})\vec{u_E})   $\\
    v\_adv\_E\_rest\_gf\_tt &$ \mu_i\nc(\frac{N_i\vec{\nabla_\perp}\phi}{B^2}\cdot\nabla \vec{u_E})  $\\
    v\_adv\_D\_main\_gf\_tt &$  \mu_i\tau_i \nc((\Delta_\perp N_i)\vec{u_E})   $\\
    v\_adv\_D\_rest\_gf\_tt &$  \mu_i\tau_i \nc(\vec{\nabla_\perp} N_i\cdot\nabla \vec{u_E})  $\\
    v\_M\_em\_gf\_tt &$  -\nc\nc(\vec{M^{em}}\vec{b_\perp})$ REMOVED! we have no $\Gamma_1 A_\parallel$ terms \\
    v\_J\_mag\_tt &$ \nc(j_{mag,||}\vec{b_\perp})=(\tau_i/2)\nc\left(\Delta_\perp( N_i U_{i,\parallel})\vec{b_\perp}\right)  $\\
    v\_J\_bperp\_tt &$ \nc(j_{||}\vec{b_\perp}) = \nc ((N_i U_{\parallel,i} - n_e u_{\parallel,e})\vec b_\perp)$ yields the Maxwell stress\\
    v\_J\_D\_gf\_tt &$ \nc\left(\frac{\bhat\times(\tau_i\vec{\nabla_\perp} N_i-\tau_e\vec{\nabla_\perp} n_e)}{B}\right) $ REMOVED! use divcurvn*\\
    v\_J\_par\_tt &$ \nc((N_i U_{i,\parallel} - n_e u_{e,\parallel})\bhat)  $ REMOVED! use divjnipar-divjnepar\\
    v\_J\_NUK\_gf\_tt &$ \nc((\mu_iN_iU_{i,\parallel}^2-\mu_en_eu_{e,\parallel}^2)\KK)  $ REMOVED! use divcurvkappa*\\
    v\_J\_JAK\_tt &$ \nc(J_\parallel A_{1,\parallel}\KK)  $ REMOVED! part of v\_J\_bperp\\
    v\_S\_E\_tt &$ \mu_i\nc(\frac{S_N\np\phi}{B^2})  $\\
    v\_S\_D\_tt &$ \mu_i\tau_i\nc(\np S_{N})  $\\
    RFB\_E\_r\_GradPsip\_tt &$ -\vec{\nabla}\phi\cn\psi_p  $\\
    RFB\_GradPi\_GradPsip\_tt &$ \tau_i\vec{\nabla} n_e/n_e\cn\psi_p  $\\
\bottomrule
\end{longtable}
Some of these definitions have "\_gf" in their definitions, as they are defined with the gyro-fluid ion density $N_i$. These quantities are also saved without this label, when using the electron fluid density $n_e$ instead of $N_i$. At the same time, all the gradients and divergences in the perpendicular direction are implemented with the toroidal field approximation.


\subsection{Parallel momentum balance, "Parallel-momentum"}
The ion parallel momentum balance reads
\begin{align}
    \mu \frac{\partial}{\partial t} \left(N U_\parallel\right) &+ \mu \nc \left( NU_\parallel \left(
    \vec u_E + \tau \vec K + \mu U_\parallel^2 \KK + U_\parallel\left(\bhat + {\vec b}_\perp\right)
    \right)\right)  \nonumber \\
    &+ 2\mu \nc ( NU_\parallel \tau \KK)
    -\mu NU_\parallel\nc \tau \KK
    + \mu NU_\parallel\mathcal K_{\vn\times\bhat}(\psi) \nonumber\\
    =& -\tau \left(\bhat + {\vec b}_\perp\right)\cn N
    -N \left( \left(\bhat+{\vec b}_\perp\right)\cn \psi + \frac{\partial A_\parallel}{\partial t}\right)
    - \eta n_e(N_iU_{\parallel,i}-n_eu_{\parallel,e})
    \nonumber\\
    &+ \mu N\Lambda_U + \mu U_\parallel \Lambda_N
\end{align}
The flux surface average over the parallel momentum equations under species summation  yields up to order $\mathcal O(\delta^3)$ in the drift-ordering
\begin{align}
  \frac{\partial}{\partial t}\RA{\mu_iN_iU_{\parallel,i} }
    % \nonumber\\
    + \frac{\partial}{\partial v} \frac{\d v}{\d\psi_p} \RA{\mu_iN_iU_{\parallel,i} \frac{\bhat\times\vn\phi}{B}\cn\psi_p + \sum_s (z_s\tau_sN_s + z_s\mu_s N_sU_{\parallel,s}^2) b_{\perp}^{\;v}  }
    \nonumber\\
   = \sum_s\RA{-z_s\tau_s \npar N_s}
   \label{eq:parallel_momentum}
\end{align}
while the toroidal parallel angular momentum contribution reads up to order $\mathcal O(\delta^3)$
\begin{align}\label{eq:parallel_momentum_direction}
    \frac{\partial}{\partial t}  \RA{\mu_iN_iU_{\parallel,i} b_\varphi}
    + \frac{\partial}{\partial v} \frac{\d v}{\d\psi_p} \RA{\mu_iN_iU_{\parallel,i} b_\varphi\frac{\bhat\times\vn\phi}{B}\cn\psi_p + \sum_s (z_s\tau_s N_s + z_s\mu_sN_sU_{\parallel,s}^2) b_\varphi b_{\perp}^{\;v} }
    \nonumber\\
    = R_0 \RA{F_{L,\varphi}}
\end{align}
The $R_0$ in the term on the right hand side stems from the normalisation of $\psi_p$.
\begin{tcolorbox}[title=Note]
The source terms add up to zero with our choice of the parallel velocity source~\eqref{eq:velocity_source}.
\end{tcolorbox}
The relevant terms in the output file are (the Lorentz force term is described in the previous subsection \ref{sec:vorticity_eq})
\begin{longtable}{llll}
\toprule
\rowcolor{gray!50}\textbf{Name} &  \textbf{Equation} &
\textbf{Name} &  \textbf{Equation}\\
\midrule
    niui &$\mu_i N_i U_{\parallel,i}$ &
    niuibphi &$\mu_i N_iU_{\parallel,i}b_\varphi$ \\
    jsparexbi\_tt       & $\mu_i N_iU_{\parallel,i}(\bhat\times\vn\phi)\cn \psi_p/B$ &
    jsparbphiexbi\_tt   & $\mu_i N_iU_{\parallel,i}b_\varphi(\bhat\times\vn\phi)\cn \psi_p/B$ \\
    divparexbi\_tt       & $\mu_i \nc (N_iU_{\parallel,i}\vec u_E)$ &
    divparbphiexbi\_tt   & $\mu_i\nc( N_iU_{\parallel,i}b_\varphi \vec u_E)$ \\
    divparexbii\_tt       & $\mu_i \nc (N_iU_{\parallel,i}\vec u_E^i)$ &
    divparbphiexbii\_tt   & $\mu_i\nc( N_iU_{\parallel,i}b_\varphi \vec u_E^i)$ \\
    divpardiai\_tt       & $\nc(\mu_i \tau_i N_iU_{\parallel,i}\vec K)$ &
    divparbphdiai\_tt   & $\nc ( \mu_i \tau_i N_iU_{\parallel,i}b_\varphi\vec K)$ \\
    divparkappai\_tt       & $\nc (\mu_i N_iU_{\parallel,i} ( \mu_i U_{\parallel,i}^2 + 2\tau_i)\KK )$ &
    divparbphikappai\_tt       & $\nc( \mu_i N_iU_{\parallel,i}b_\varphi ( \mu_i U_{\parallel,i}^2 + 2\tau_i)\KK)$ \\
    divparmirrorAi\_tt       & $\nc(z_i \tau_i N_i \vec b_\perp)$ &
    divparbphiApar\_tt   & $\sum_s \nc ((z_s \tau_s N_s + z_s \mu_s N_s U_{\parallel,s}^2)b_\varphi \vec b_\perp)$ \\
    divparmirrorAe\_tt       & $-\nc ( n_e \vec b_\perp)$ &
    & \\
    divparApari\_tt       & $\nc( z_i \mu_i N_i U_{\parallel,i}^2\vec b_\perp)$ &
    & \\
    divparApare\_tt       & $\nc( z_e \mu_e n_e u_{\parallel,e}^2\vec b_\perp)$ &
    & \\
    divjpari\_tt & $ \nc( +\mu_i N_i U_{\parallel,i}^2 \bhat)$ &
    & \\
    divjpare\_tt & $ \nc( -\mu_e n_e u_{\parallel,e}^2 \bhat)$ &
    & \\
    lparpar\_tt   & $N_i \Lambda_{U_i,\parallel} + U_{\parallel,i} \Lambda_{N_i,\parallel}$ &
    lparperp\_tt & $U_{\parallel,i} \Lambda_{N_i} + N_i\Lambda_{U_i} $ \\
    lparparbphi\_tt   & $(N_i \Lambda_{U_i,\parallel} + U_{\parallel,i} \Lambda_{N_i,\parallel})b_\varphi$ &
    lparperpbphi\_tt & $(U_{\parallel,i} \Lambda_{N_i} + N_i\Lambda_{U_i})b_\varphi $ \\
    sparKappaphii\_tt & $-\mu_i N_i U_{\parallel,i} \KK\cn \psi_i$ &
    sparmirrorKappai\_tt & $\tau_i N_i U_{\parallel,i} \nc\KK$ \\
    sparphii\_tt & $-N_i\npar \psi$ &
    friction\_tt & $ \eta n_e(N_iU_{\parallel,i}-n_eu_{\parallel,e})$ \\
    sparmirrori\_tt & $-\tau_i\npar N_i$ &
    sparmirrorAi\_tt & $-\tau_i \vec b_\perp \cn N_i $ \\
    sparphiAi\_tt & $ -N_i \vec b_\perp \cn \psi $ &
    spardotAi\_tt & $ -N_i \partial A_\parallel /\partial t$ \\
\bottomrule
\end{longtable}
Note that the parallel viscosity term vanishes exactly under the flux-surface average. This can serve as a numerical test.

We gather the dominant terms in the electron momentum equation (neglecting all terms as $\mu_e=0$). This leaves the parallel force balance
\begin{align}
    -(\bhat + \vec b_\perp) \cn n_e +n_e\left( \left( \bhat + \vec b_\perp \right) \cn \phi + \frac{\partial A_\parallel}{\partial t} \right) +\eta n_e \left( N_iU_{\parallel,i} - n_eu_{\parallel,e}\right) = 0
\end{align}
Further we can compare if the following terms are small in the electron force balance
\begin{align*}
    \nc( z_e \mu_e n_e u_{\parallel,e}^2\vec b_\perp)\\
    \RA{\nc ( \vec b_\perp n_e )} = \RA{n_e \vec b_\perp \cn n_e}
\end{align*}
\begin{longtable}{llll}
\toprule
\rowcolor{gray!50}\textbf{Name} &  \textbf{Equation} &
\textbf{Name} &  \textbf{Equation}\\
\midrule
    sparphie\_tt & $n_e\npar \phi$ &
    friction\_tt & $ \eta n_e(N_iU_{\parallel,i}-n_eu_{\parallel,e})$ \\
    sparmirrore\_tt & $-\npar n_e$ &
    sparmirrorAe\_tt & $-\vec b_\perp \cn n_e $ \\
    sparphiAe\_tt & $n_e \vec b_\perp \cn \phi$ &
    spardotAe\_tt & $ n_e \partial A_\parallel /\partial t$ \\
    neue &$n_e u_{\parallel,e}$ &
    & \\
\bottomrule
\end{longtable}


%%%%%%%%%%%%%%%%%%%%%%%%%%%%%%%%%%%%%%%


\subsection{Zonal flow energy, "Zonal-Flow-Energy"}

\begin{align}
    E_\mathrm{zonal} = \frac{1}{2}\RA{\rho_M}\FA{ \iota^{-2}\mathcal I^{\vartheta\vartheta} + 2\iota^{-1}\mathcal I^{\vartheta\varphi} + \mathcal I^{\varphi\varphi}} \FA{u_{E,\varphi}}^2
    \equiv \frac{1}{2} \RA{\rho_M} \FA{u_{E,\varphi}}^2  \FA{\mathcal I_0}
    \label{eq:zonal_energy}
\end{align}
For symmetry flux coordinates we have $g_{\vartheta\vartheta} = R^2 (\vn\psi_p)^2/I^2\iota^2$, $g_{\varphi\vartheta} =0$ and $g_{\varphi\varphi}=R^2$ and thus $\mathcal I_0 = R^{-2}( 1 + I^2/|\vn\psi_p|^2)= B^2 / (R_0^2|\vn\psi_p|^2)$. (In dimensionless notations; note that $g_{ij}$ has units of $\rho_s^2$ and $\mathcal I_0$ has units $\rho_s^{-2}$)
\begin{align}\label{eq:perp_kinetic}
      \frac{\partial}{\partial t}E_{\mathrm{zonal}} +\frac{\partial}{\partial v } \left(E_{\mathrm{zonal}}\FA{u^v} \right)
  =&-\FA{\mathcal I_0}\FA{u_{E,\varphi}}\left(\frac{\partial }{\partial v}  \Theta_{\varphi}^{\; v} + \RA{(\vec j_f\times\vec B)_\varphi}\right)
  \nonumber\\
    &-\frac{1}{2}\FA{u_{E,\varphi}}^2\frac{\partial}{\partial v}\left(\RA{\rho_M}\FA{\FF{\mathcal I_0}\FF{u^v}}\right)
     + \mathcal S_{\mathrm{zonal}}
\end{align}
where we neglected the term $\RA{n\vec u\cn \mathcal I_0}$ in the continuity equation as small in our ordering
 and we have
 \begin{align}
 \mathcal S_{\mathrm{zonal}} :=& \FA{\mathcal I_0} \FA{u_{E,\varphi}} \mathcal S_{u_{E,\varphi}} + \frac{1}{2}\FA{u_{E,\varphi}}^2  \RA{mS_n\mathcal I_0}
%  \nonumber\\
 \label{eq:zonal_source}
 \end{align}
 and in the output file
\begin{longtable}{llll}
\toprule
\rowcolor{gray!50}\textbf{Name} &  \textbf{Equation} &
\textbf{Name} &  \textbf{Equation}\\
\midrule
    nei0 &$n_e \mathcal I_0$ &
    snei0\_tt & $S_{n_e } \mathcal I_0$ \\
\bottomrule
\end{longtable}

%%%%%%%%%%%%%%%%%%%%%%%%%%%%%%%%%%%%%%%%%%%%%%%%%%%%%%%%%%%%%%%%%%%%%%%%%
\section{Probes Output}\label{sec:probes-output}
The probes output is stored in a netcdf group called "probes" in the main output file
%
%Name | Type | Dimensionality | Description
%---|---|---|---|
\begin{longtable}{lll>{\RaggedRight}p{7cm}}
\toprule
\rowcolor{gray!50}\textbf{Name} &  \textbf{Type} & \textbf{Dimension} & \textbf{Description}  \\ \midrule
format           & text attribute   & - & The format value as a styled string \\
dim              & Coord. Var. & 1 (dim) & index-like-coordinate $\in [0,1]$ \\
time             & Coord. Var. & 1 (time)& time at which fields are written (variable size: itstp*maxout$+1$, dimension size: unlimited) \\
X                & Dataset & 1 (dim)      & Static field data \\
Y                & Dataset & 2 (time,dim) & Dynamic measurements \\
\bottomrule
\end{longtable}
where $X\in $["R","Z","P", "a collection of magnetic field quantities (check output file for details"]
and $Y\in $
\begin{longtable}{llllll}
\toprule
\rowcolor{gray!50}\textbf{Name} &  \textbf{Equation} & \textbf{Name} &  \textbf{Equation} & \textbf{Name} & \textbf{Equation}\\
\midrule
    ne &$n_e$ &
    ni &$N_i$ &
    ue &$u_{\parallel,e}$ \\
    ui &$U_{\parallel,i}$ &
    phi &$\phi$ &
    apar &$A_\parallel$ \\
    neR &$\partial_R n_e$ &
    niR &$\partial_R N_i$ &
    ueR &$\partial_R u_{\parallel,e}$ \\
    uiR &$\partial_R U_{\parallel,i}$ &
    phiR &$\partial_R \phi$ &
    aparR &$\partial_R A_\parallel$ \\
    neZ &$\partial_Z n_e$ &
    niZ &$\partial_Z N_i$ &
    ueZ &$\partial_Z u_{\parallel,e}$ \\
    uiZ &$\partial_Z U_{\parallel,i}$ &
    phiZ &$\partial_Z \phi$ &
    aparZ &$\partial_Z A_\parallel$ \\
    nePar &$\npar n_e$ &
    niPar &$\npar N_i$ &
    uePar &$\npar u_{\parallel,e}$ \\
    uiPar &$\npar U_{\parallel,i}$ &
    phiPar &$\npar \phi$ &
     & \\
\bottomrule
\end{longtable}
with all these quantities available most other quantities (flux, energy, etc.)
can be computed in post processing.

%%%%%%%%%%%%%%%%%%%%%%%%%%%%%%%%%%%%%%%%%%%%%%%%%%%%%%%%%%%%%%%%%%%%%%%%%
\section{Diagnostics}\label{sec:diagnostics}
\mintinline{bash}{feltor/src/feltor/feltordiag.cu}
reads one or more previously generated simulation file(s) \mintinline{bash}{input0.nc ... inputN.nc}
described in Section~\ref{sec:output_file} and writes into a
 single second output file \mintinline{bash}{output.nc} described as follows. \\
Compilation
\mint{bash}{make feltordiag device={gpu,omp}}
\noindent Usage
\mint{bash}{./feltordiag configuration.json input0.nc ... inputN.nc output.nc}
The usage requires a file \mintinline{bash}{configuration.json} where parameters for the diagnostic can be introduced and is described further down.
\begin{tcolorbox}[title=Note]
\mintinline{bash}{feltordiag} refuses to overwrite existing files in order to protect against data loss in case of accidental spelling
errors or other careless mistakes.
\end{tcolorbox}

Output file format: \href{https://www.unidata.ucar.edu/software/netcdf/docs/}{netcdf-4/hdf5};
\href{http://cfconventions.org/Data/cf-conventions/cf-conventions-1.7/cf-conventions.html}{CF Conventions CF-1.7}

A \textit{coordinate variable (Coord. Var.)} is a Dataset with the same name as a dimension.

\begin{longtable}{lll>{\RaggedRight}p{7cm}}
\toprule
\rowcolor{gray!50}\textbf{Name} &  \textbf{Type} & \textbf{Dimension} & \textbf{Description}  \\ \midrule
inputfile  &     text attribute & - & verbose input file as a string (valid JSON, C-style comments are allowed but discarded) \\
configfile   &     text attribute & - & verbose configuration file as a string (valid JSON, C-style comments are allowed but discarded) \\
x                & Coord. Var. & 1 (x) & $R$-coordinate (computational space, compressed size: $nN_x/c_x$)\\
y                & Coord. Var. & 1 (y) & $Z$-coordinate (computational space, compressed size: $nN_y/c_y$)\\
psi              & Coord. Var. & 1 (psi) & $\psi_p$-coordinate ( default size: $3\cdot 64$) \\
eta              & Coord. Var. & 1 (eta) & $\eta$-coordinate ( default size: $3\cdot 640$) \\
xc               & Dataset & 2 (eta,psi) & Cartesian x-coordinate of the FSA grid \\
yc               & Dataset & 2 (eta,psi) & Cartesian y-coordinate of the FSA grid\\
vol              & Dataset & 2 (eta,psi) & Volume form in 2d including dG weights (can be used to integrate 2d quantities on the FSA grid; multiply by $2\pi$xc to get 3d volume)  \\
time             & Coord. Var. & 1 (time)& time at which fields are written (variable size: maxout$+1$, dimension size: unlimited) \\
dvdpsip          & Dataset & 1 (psi) & $\d v/\d\psi_p$ \\
psi\_vol         & Dataset & 1 (psi) & The volume enclosed by the flux surfaces $v(\psi_p) = \int_{\psi_p} \dV $ \\
psi\_area        & Dataset & 1 (psi) & The area of the flux surfaces $A(\psi_p) = 2\pi \int_\Omega |\vn\psi_p| \delta(\psi_p - \psi_{p0}) H(Z-Z_X) R\d R\d Z$ \\
dv2ddpsip        & Dataset & 1 (psi) & $\d v/\d\psi_p$ (in the 2d poloidal plane) \\
psi\_vol2d       & Dataset & 1 (psi) & The 2d volume enclosed by the flux surfaces $v(\psi_p) = \int_{\psi_p} dRdZ $ \\
psi\_arc         & Dataset & 1 (psi) & The arc length of the flux surfaces $L(\psi_p) = 2\pi \int_\Omega |\vn\psi_p| \delta(\psi_p - \psi_{p0}) H(Z-Z_X) \d R\d Z$ \\
q-profile        & Dataset & 1 (psi) & The safety factor $q(\psi_p)$ \eqref{eq:safety_factor} using direct integration ( accurate but we assign random values outside separatrix) \\
psi\_psi         & Dataset & 1 (psi) & explicit $\psi_p$ values; Same as psi \\
psit1d           & Dataset & 1 (psi) & Toroidal flux (integrated q-profile) $\psi_t = \int^{\psi_p} \d\psi_p q(\psi_p)$ \\
rho              & Dataset & 1 (psi) & Transformed flux label $\rho:= 1 - \psi_p/\psi_{p,O}$ \\
rho\_p           & Dataset & 1 (psi) & poloidal flux label $\rho_p:= \sqrt{1 - \psi_p/\psi_{p,O}}$  ( see Section~\ref{sec:alternative}\\
    rho\_t           & Dataset & 1 (psi) & Toroidal flux label $\rho_t :=
    \sqrt{\psi_t/\psi_{t,\mathrm{sep}}}$ (is similar to $\rho$ in the edge but
    $\rho_t$ is nicer in the core domain, because equidistant $\rho_t$ make
more equidistant flux-surfaces, see Section~\ref{sec:alternative} )\\
Z\_fluc2d        & Dataset & 3 (time,y,x) & Fluctuation level on selected plane ($\varphi= 0$) $\delta Z := Z(R,Z,0) - \RA{ Z}(R,Z)$ \\
Z\_fsa2d         & Dataset & 3 (time, y,x) & Flux surface average $\RA{ Z}$ interpolated onto 2d plane Eq.~\eqref{eq:fsa_vol} \\
Z\_cta2d         & Dataset & 3 (time, y,x) & Convoluted toroidal average Eq.~\eqref{eq:cta} \\
Z\_fsa           & Dataset & 2 (time, psi) & Flux surface average $\RA{ Z}$ Eq.~\eqref{eq:fsa_vol} \\
Z\_std\_fsa      & Dataset & 2 (time, psi) & Standard deviation of flux surface average on selected plane ($\varphi=0$) $\sqrt{\RA{(\delta Z)^2}}$ \\
Z\_ifs           & Dataset & 2 (time, psi) & Volume integrated flux surface average $\int\d v\RA{ Z}$ unless Z is a current, then it is the volume derived flux-surface average $\partial_v \RA{ Z}$ \\
Z\_ifs\_lcfs     & Dataset & 1 (time) & Volume integrated flux surface average evaluated on last closed flux surface $\int_0^{v(0)}\d v\RA{ Z}$ unless Z is a current, then it is the fsa evaluated $\RA{ j_v}(0)$ \\
Z\_ifs\_norm     & Dataset & 1 (time) & Volume integrated square flux surface average $\sqrt{\int \d v \RA{Z}^2}$, unless Z is a current, then it is the square derivative of the flux surface average $\sqrt{\int\d v (\partial_v \RA{j^v})^2}$\\
Z\_cta2dX  & Dataset & 3 (time, psi, eta) & Convoluted toroidal average Eq.~\eqref{eq:cta} in the magnetic plane. \\
\bottomrule
\label{table:diag}
\end{longtable}
where Z $\in$ \{X, Y\_tt\}
\begin{tcolorbox}[title=Note]
    \mintinline{bash}{feltoridag} converts all $jsX$ quantities into $jvX$
by multiplying $\d v/\d \psi_p$
in the sense that $\vec j\cn v  = \vec j \cn \psi_p \d v/\d\psi_p$.
\end{tcolorbox}
The default parameters used for the X-point
\href{https://feltor-dev.github.io/doc/geometries/html/structdg_1_1geo_1_1_separatrix_orthogonal.html}{separatrix-orthogonal Grid}
construction are $f_x = 1/8$, $f_y = 0$, $n_\psi = 3$, $N_\zeta = 64$ and
$N_\eta = 640$ and the constant monitor metric~\cite{Wiesenberger2018}, but can be specified in the \mintinline{bash}{configuration.json} file.
Furthermore, a user can specify which kind of diagnostics quantities should be computed and how:
\begin{minted}[texcomments]{js}
// for feltordiag
"n": 3, // resolution of X-point grid for fsa
"Npsi": 64, // resolution of X-point grid for fsa
"Neta": 640, // resolution of X-point grid for fsa
"Kphi": 10, // resolution for convoluted toroidal average (10 is fine)
"fsa" : "toroidal-average", // From which quantity fsa should be computed, can also be "convoluted-toroidal-average"
"fx_0" : 0.125, // where the separatrix is in relation to the $\zeta$ coordinate
"x-grid-interpolation" : "dg", // type of interpolation for Cylindrical2XGrid
"cta-interpolation" : "dg", // interpolation type in Fieldaligned object for convoluted toroidal averages
"diagnostics" :
[
    "fsa",  // save \_fsa in table \ref{table:diag}.
    "fsa2d", //  save \_fsa2d in table \ref{table:diag}.
    "cta2d", //  save \_cta in table \ref{table:diag}.
    "cta2dX", //  save \_cta2dX in table \ref{table:diag}.
    "fluc2d", //  save \_fluc2d in table \ref{table:diag}.
    "ifs", //  save \_ifs in table \ref{table:diag}.
    "std_fsa", //  save \_std\_fsa in table \ref{table:diag}.
    "ifs_lcfs", //  save \_ifs\_lcfs in table \ref{table:diag}.
    "ifs_norm" //  save \_ifs\_norm in table \ref{table:diag}.
],
// for interpolate in 3d
"fine-grid-factor" : 12, // how many planes are outputted in the new file in total: Nz*fine-grid-factor
"time-reduction-factor" : 10 // only ouptut every *th step
\end{minted}
\begin{tcolorbox}[title=Note]
    If the "diagnostics" field does not exist, all quantities will be selected and written to the output file
\end{tcolorbox}


%%%%%%%%%%%%%%%%%%%%%%%%%%%%%%%%%%%%%%%%%%%%%%%%%%%%%%%%%%%%%%%%%%%%%%%%%%%%%%%%%%%%
\subsection{Flux surface averaging and safety factor}
\subsubsection{Preliminary}
Recall that the {\bf Dirac delta-function} has the property (in any dimension):
\begin{align} \label{eq:dirac_delta}
\int_V f(\vec x) \delta(h(\vec x) - h') \dV = \int_{h=h'} \frac{f(\vec x)}{|\vn h|} \dA
\end{align}
which means that the delta-function can be used to express area integrals of the
submanifold given as a contour of the function $h(\vec x)$.
A numerically tractable approximation to the delta-function reads
\begin{align}\label{eq:delta}
\delta(h(\vec x)-h') = \frac{1}{2\pi \epsilon^2}
\exp\left( - \frac{\left(h(\vec x)-h'\right)^2}{2\epsilon^2}\right)
\end{align}
where $\epsilon$ is a small, free parameter.
In the DG framework the left-hand side
of Eq.~\eqref{eq:dirac_delta} can thus readily be computed
via Gauss-Legendre quadrature, which we propse as a first method to compute area
integrals even if our coordinate system is not aligned to the area.
Note: in order for this to work the Delta function needs to be numerically
resolved and cannot be made arbitrarily small.
This introduces a smoothing effect
over neighboring contour lines which is given by the grid distance.

Furthermore, recall the {\bf co-area formula}
\begin{align} \label{eq:coarea}
\int_{\Omega_0} f(\vec x) \dV =
\int_0^{h_0} \left( \int_{h=h'} \frac{f(\vec x)}{|\vn h|}  \dA  \right) \d h'
\end{align}
where $\Omega_0$ is the volume enclosed by the contour $h=h_0$.
The co-area formula can be viewed as a change of variables in the
volume integral.

We define the {\bf toroidal average} of a function $f(R,Z,\varphi)$ as
\begin{align} \label{eq:phi_average}
\PA{ f}(R,Z) := \frac{1}{2\pi}\oint f(R,Z,\varphi)\d \varphi
\end{align}

In arbitrary coordinates the area 2-form is given by the interior product of the
surface normal with the volume form
\begin{align}
\label{}
\dA^2 = i_{\hat \psi_p} vol^3 \quad \hat \psi_p = \frac{\vn \psi_p}{|\vn \psi_p|}
\end{align}
The pullback to a flux-aligned coordinate system $\{\zeta, \eta, \varphi\}$ is trivial (since $\vn \psi_p$ only has one component $g^{\zeta\zeta} \partial \psi_p / \partial \zeta$) and we have
\begin{align}
\dA &= \sqrt{g^{\zeta\zeta}} \sqrt{g} \d\eta\d\varphi = f_0|\vn\psi_p|\sqrt{g}\d\eta\d\varphi,
\\
\vec\dA &:= \hat\psi_p \dA = f_0 (\vn\psi_p) \sqrt{g}\d\eta\d\varphi,\quad
\label{}
\end{align}
where we used that $g^{\zeta\zeta} = (\vn\zeta)^2 = f_0^2(\vn\psi_p)^2$.
Notice that numerically we can integrate in flux-aligned coordinates by generating a corresponding
grid and pulling back (interpolating) the relevant fields to this grid. This is the second method
to numerically compute area integrals.

%%%%%%%%%%%%%%%%%%%%%%%%%%%%%%%%%%%%
\subsubsection{Flux surface average}


The flux surface average (as a {\bf volume average} after \cite{haeseleer}) is defined as an average over a
small volume - a shell centered around the flux-surface - defined by two neighboring flux-surfaces.
With the help of the volume
flux label (notice that both the volume $v$ as well as the poloidal flux $\psi_p$ have physical
meaning while the coordinate $\zeta(\psi_p)$ is an arbitrary choice) we define
\begin{align} \label{eq:fsa_vol}
v(\psi_p) :=& \int_{\psi_{p,O}}^\psi \dV = \int^{\zeta(\psi_p)} \sqrt{g}\d\zeta\d\eta\d\varphi,
\\
\frac{\d v}{\d\psi_p} =& \int\dA |\vn\psi_p|^{-1} = 2\pi f_0\oint_{\zeta(\psi_p)} \sqrt{g}\d\eta \\
\RA{ f }_\psi :=& \frac{\partial}{\partial v} \int \dV f
 = \frac{1}{\int \dA |\vn\psi_p|^{-1} } \int_{\psi_p} \frac{f(\vec x)}{|\vn\psi_p|} \dA \nonumber\\
=& \frac{\int_\Omega \PA{ f}(R,Z) \delta(\psi_p(R,Z)-\psi_{p})H(Z-Z_X)\ R \d R \d Z}
{\int_\Omega \delta(\psi_p(R,Z)-\psi_{p})H(Z-Z_X)\ R \d R \d Z}\nonumber\\
 =& \left(\frac{\d v}{\d\psi_p }\right)^{-1} 2\pi f_0 \oint_0^{2\pi} \PA{ f}(\zeta,\eta) \sqrt{g}\d\eta
 = \frac{1}{\oint \sqrt{g}\d\eta } \oint_0^{2\pi} \PA{ f}(\zeta,\eta) \sqrt{g}\d\eta
\end{align}
where we used the co-area formula Eq.~\eqref{eq:coarea} for the second identity
and we use the Heaviside function $H(Z-Z_X)$ to cut away contributions from below the X-point
in our domain $\Omega$.
 We immediately see that this definition is particularly easy to compute
 in a flux-aligned coordinate system. Notice however that the volume element
 does appear (unlike e.g. Tokam3X papers).
 We use our grid construction algorithm with constant monitor metric described in Reference~\cite{Wiesenberger2018} to construct a flux-aligned grid and interpolate
 the values of any function onto its grid points.
 Even though this grid is unusable for simulations due to the diverging metric at the X-point the
 evaluation of integrals works well as the singularity is integrable.

The flux-surface average fulfills the basic identities
\begin{align}
\label{eq:fsa_identities}
\RA{ \mu f + \lambda g} &= \mu\RA{ f} + \lambda \RA{ g} \\
\RA{ f(\psi_p)} &= f(\psi_p)
\end{align}

The volume average is well-suited for density-like quantities
as we can see with the following identity.
Assume we have a quantity $X$ with $\partial_t X + \nc \vec j_X = \Lambda_X$.
Then we can use the volume average to write
\begin{align}
\frac{\partial}{\partial t} \RA{X } + \frac{\partial}{
  \partial v} \RA{ \vec j_X\cn v}  = \RA{ \Lambda_X}
\label{eq:fsa_balance}
\end{align}
where again $v=v(\psi_p)$ is the volume flux label.
The {\bf total flux} of a given flux density $\vec j_X$ through the
flux surface $\psi_p = \psi_{p0}$ is given by
\begin{align}
\RA{\vec j_X\cn v} &:= J_X=\oint_{\psi_p=\psi_{p0}} \vec j_X\cdot \vec{\dA} =
 \frac{\d v}{\d\psi_p} \RA{ \vec j_X\cn\psi_p }\nonumber\\
 &=
   2\pi f_0 \oint_0^{2\pi} \PA{ \vec j_X\cn\psi_p}(\zeta,\eta) \sqrt{g}\d\eta
%2\pi\int_\Omega \vec \PA{ \vec j\cn\psi_p} \delta(\psi_p(R,Z)-\psi_{p0}) H(Z-Z_X)\ R \d R \d Z
\label{eq:total_flux}
\end{align}
Once we have the flux-surface averaged equation we can easily get the volume integrated version (again with the help of the co-area formula)
\begin{align}
\frac{\partial}{\partial t} \int_0^{v(\psi_p)}\RA{X} \d v 
+ \RA{ \vec j_X\cn v}(v(\psi_p))  = \int_0^{v(\psi_p)}\RA{ \Lambda_X}\d v
\label{eq:integral_balance}
\end{align}

\subsubsection{The safety factor}
Assume that we pick a random field line and follow it (integrate it) for exactly one
poloidal turn. The {\bf safety factor} is defined as the ratio between
the resulting toroidal angle ($\Delta\varphi$) to the poloidal angle ($2\pi$)
\begin{align}
q := \frac{\Delta\varphi}{2\pi}
\label{}
\end{align}
Since our magnetic field is symmetric in $\varphi$ and we used one
full poloidal turn this definition is independent of which
fieldline we pick on a given flux surface.

%We define the poloidal length $s$ as the fieldline following
%parameter i.e. $\vec B\cn s \equiv B_p = R_0|\vn \psi_p|/R$
%and $\d\varphi/\d s = B^\varphi(R(s), Z(s)) / B_p(R(s),Z(s))$.
%We can then express the safety factor as the line integral
%\begin{align}
%q=\frac{1}{2\pi}\oint \frac{B^\varphi}{B_p} \d s = \frac{1}{2\pi}\oint_{\psi_p=\psi_{p0}}\frac{I(\psi_p)}{R|\vn\psi_p|} \d s
%= \frac{1}{2\pi}\int \frac{I(\psi_p)}{R}\delta(\psi_p-\psi_{p0}) H(Z-Z_X) \d R\d Z
%\end{align}
%where we made use of Eq.~\eqref{eq:dirac_delta} in two dimensions in the
%last equality and thus arrive at a numerical tractable expression
%to evaluate the safety factor.
We define the geometric poloidal angle $\Theta$ as the fieldline following
parameter i.e. $\vec B\cn\Theta = R_0(\psi_R (R-R_0) + \psi_Z Z)/r^2R$.
We can then directly integrate the safety factor as
\begin{align}\label{eq:safety_factor}
\frac{\d R}{\d\Theta} = \frac{B^R}{B^\Theta}\quad
\frac{\d Z}{\d\Theta} = \frac{B^Z}{B^\Theta}\quad
\frac{\d \varphi}{\d\Theta} = \frac{B^\varphi}{B^\Theta}\\
q\equiv\frac{1}{2\pi}\oint \frac{B^\varphi}{B^\Theta} \d\Theta
\end{align}
We integrate this equation with the help of one of our ODE integrators, i.e. we use a high-order Runge-Kutta method
and refine the stepsize until machine-precision is reached.
\begin{tcolorbox}[title=Note]
The safety factor diverges on the last closed flux
surface whereas Eq.~\eqref{eq:total_flux}
remains finite due to the $\vn\psi_p$ factor. Outside the LCFS $q$ remains undefined.
\end{tcolorbox}
%%%%%%%%%%%%%%%%%%%%%%%%%%%%%%%%%%%%%%%%%%%%%%%%%%%%%%%%%%%%%%%%%%%%%%%%%%%%%%%%%%%%%%%%%%%%
\subsubsection{Toroidal averages}
Here, we comment on the $\varphi$ average that is part of the flux-surface average Eq.~\eqref{eq:fsa_vol}.
One simple approach is
quadrature of the form
\begin{align}\label{eq:toroidal_summation}
    \bar f = \frac{1}{N} \sum_{i=0}^{N-1} f_i (R,Z)
\end{align}
where $N=32$ in most of our simulations and $f_i$ is the $i$-th toroidal plane.
Since the boundary conditions in $\varphi$ are periodic this amounts to the trapezoidal rule.
A low number of toroidal planes is sufficient in simulations when we use the toroidal field
approximation in combination with the
flux-coordinate independent (FCI) approach for the parallel derivatives.
However, since the actual $\varphi$ direction is
under-resolved the
integration gives a wrong answer to the actual $\varphi$ average (seen in 2d plots as little humps).
This is because the resulting structures are predominantly field aligned and not toroidally symmetric.

In order to improve the toroidal average we now have the following idea:
if we, before we do the $\varphi$ integration,
interpolate the function to integrate onto a large number of toroidal
planes then the result should
be more accurate than before.
In other words we interpolate the function given on the coarse $\varphi$ simulation grid
onto a hypothetic fine $\varphi$ grid along the magnetic field lines
and only then compute the $\varphi$ average.

Let us divide the $\varphi$ direction between two original planes into $N_\varphi+1$ (a large number) equidistant planes
of distance $\delta \varphi$ and integrate the magnetic field $\vec B$ in between.
\begin{subequations}
\begin{align}
    \frac{\d R}{\d\varphi}&= \frac{B^R}{B^\varphi},\\ %\frac{R}{I}\frac{\partial\psi}{\partial Z},\\
    \frac{\d Z}{\d\varphi}&=\frac{B^Z}{B^\varphi},\\%-\frac{R}{I}\frac{\partial\psi}{\partial R}.
\end{align}
\label{eq:fieldline}
\end{subequations}
We integrate Eqs.~\eqref{eq:fieldline} from $\varphi=0$ to $\varphi=\pm \Delta \varphi$
with initial condition
\begin{align}
    (R(0), Z(0) ) = (R, Z).
    \label{}
\end{align}
Let us characterize the solution $(R(\pm \delta \varphi), Z(\pm \delta \varphi))$ to Eqs.~\eqref{eq:fieldline} as the flow generated by $\vec B/B^\varphi$
\begin{align}
    \Tdpm\vec z \equiv \Tdpm[R, Z, \varphi]:= ( R(\pm \delta\varphi), Z( \pm \delta\varphi), \varphi\pm\delta \varphi),
    \label{}
\end{align}
Obviously we have $\Tdm\circ\Tdp = 1$, but $\Tdpm$ is not necessarily unitary since $\vec B/B^\varphi$ is in general
not divergence free.
We are now able to extend the function $f$ given on the coarse $\varphi$ grid unto the fine $\varphi$ grid via
\begin{align}
    f(R,Z,\varphi_0+j\delta\varphi) = {\Tdm}^j f(R,Z,\varphi_0)\\
    f(R,Z,\varphi_0-j\delta \varphi) = {\Tdp}^j f(R,Z,\varphi_0)
\end{align}
This gives simple 0-th order extrapolation of our function.
Let us call $f_i := f(R,Z,\varphi_i)$ the $i$-th toroidal plane and $N_\varphi$ even. Then
we have the following integration, where we consider the original toroidal planes as cell-centered
\begin{align}
    \RA{f}_\varphi &= \frac{1}{(N_\varphi+1) N} \left[\left(
    {\Tdp}^{N_\varphi/2} f_0 + ... + \Tdp f_0 + f_0 + \Tdm f_0 ... + {\Tdm}^{N_\varphi/2} f_0\right)\right. \nonumber\\
    &\left. +\left( {\Tdp}^{N_\varphi/2} f_1 + ... + \Tdp f_1 + f_1 + \Tdm f_1 ... + {\Tdm}^{N_\varphi/2} f_1\right)  + ... \right] \nonumber\\
    &= \frac{1}{N (N_\varphi+1)} \sum_{i=0}^{N-1} \left[f_i + \sum_{j=1}^{N_\varphi/2} \left( {\Tdm}^j f_i + {\Tdp}^jf_i\right)\right] \nonumber\\
    &=
    \frac{1}{N_\varphi+1} \left[ \sum_{j=0}^{N_\varphi/2}  {\Tdm}^j \left(\frac{1}{N}\sum_{i=0}^{N} f_i\right)
    +
    \sum_{j=1}^{N_\varphi/2}  {\Tdp}^j \left(\frac{1}{N}\sum_{i=0}^{N} f_i\right)\right]
    \nonumber\\
    &=
    \frac{1}{N_\varphi+1} \left[ \sum_{j=0}^{N_\varphi/2}  {\Tdm}^j \bar f(R,Z)
    +
    \sum_{j=1}^{N_\varphi/2}  {\Tdp}^j \bar f(R,Z)\right]
\end{align}
Here, we used that the push-forward operator $\Tdm$ is linear that is $\Tdm f_0 + \Tdm f_1 = \Tdm (f_0+f_1)$
and recover the simple toroidal summation $\bar f$ Eq.~\eqref{eq:toroidal_summation}.
Now, we can see that in the limit $N_\varphi \rightarrow\infty$ the discrete sum represents the integral
of the form
\begin{align}
    \RA{f}_\varphi(R,Z) = \frac{1}{\Delta\varphi}\int_{-\Delta\varphi/2}^{\Delta\varphi/2}\d\varphi \bar f(R(\varphi),Z(\varphi))
\end{align}
A consistency test of this approach is to simply use $\vec B = e_\varphi$. Then
$\Tdpm = 1$ and we recover the original integration $\RA{f}_\varphi= \bar f$.
Now, instead of doing a 0-th order interpolation let us try a linear interpolation along field-lines in between planes that is ( assuming $N_\varphi$ toroidal planes)
\begin{align}
    f(R,Z,\varphi_i + j\delta\varphi) = \left(1-\frac{j}{N_\varphi}\right){\Tdm}^j f_i + \frac{j}{N_\varphi} {\Tdp}^{N_\varphi -j}f_{i+1}
\end{align}
\begin{align}
    \RA{ f}_\varphi &= \frac{1}{N_\varphi N} (( f_0 + (1-\alpha_1)\Tdm f_0 + \alpha_1 (\Tdp)^{N_\varphi -1} f_1 + (1-\alpha_2)(\Tdm)^2 f_0 + \alpha_2 (\Tdp)^{N_\varphi-2} f_1 ...  )+ (f_1 + ...) + ...)\nonumber\\
    &= \frac{1}{ N_\varphi} \sum_{j=0}^{N_\varphi-1}  (1-\alpha_j)(\Tdm)^j \bar f+\alpha_j (\Tdp)^{N_\varphi -j} \bar f
    = \frac{1}{\Delta\varphi}\int_{-\Delta\varphi}^{\Delta\varphi}\d\varphi w(\varphi) \bar f (R(\varphi),Z(\varphi))
    \label{eq:cta}
\end{align}
with $\alpha_j = \frac{j}{N_\varphi}$ and $w(\varphi)$ a linear weight function (pyramid shape) with $\int_{-\Delta\varphi}^{\Delta\varphi} w(\varphi) = \Delta\varphi$. Taking $\Tdpm=1$ again leads to the old result.


Now, an interesting question is, what happens if we are trying to apply the above
results to a function that is not field-aligned like $B(R,Z)$ of $\vec K\cn\psi_p$ for instance? For those functions $\bar f_\mathrm{old}$ actually yields the exact
result, while the convolution is an approximation.
Here, we have to test.
Typically, the functions that we use are slowly varying in $R$ and $Z$ and
so the convolution should not change the result too much.
A good test candidate is still $\langle \mathcal K(\psi_p)\rangle_{\psi_p}=0$.

In all practical tests so far the flux-suface average is not or only very slightly changed by this procedure.
This means that it is not
necessary to follow the smoothing procedure if one is only interested in the flux-surface average.
This makes sense because the toroidal and poloidal averages commute.

%%%%%%%%%%%%%%%%%%%%%%%%%%%%%%%%%%%%%%%%%%%%%%%%%%%%%%%%%%%%%%%%%%%%%%
\subsection{Alternative flux labels} \label{sec:alternative}
We find the toroidal flux $\psi_t$ by integrating the q-profile $\psi_t = \int^{\psi_p} \d\psi_p q(\psi_p)$. Since $q$ diverges, $\psi_t$, in contrast to $\psi_p$,
is not defined outside the last closed flux-surface (but has a finite value on the last closed flux surface). We now define the normalized poloidal and toroidal flux labels $\rho_p$ and $\rho_t$
\begin{align}
    \rho_p&:= \sqrt{1-\frac{\psi_p }{\psi_{p,O}}} \ \leftrightarrow\ \psi_p = (1-\rho_p^2)\psi_{p,O} \\
    \rho_t&:= \sqrt{\frac{\psi_t}{\psi_{t,\mathrm{sep}}}},\\
    \text{with }\psi_{p,O} &= \psi_p(R_O, Z_O)% \text{ and } \psi_{p,X} = \psi_p(R_X, Z_X)
\end{align}
where $R_O$, $Z_O$ are the coordinates of the O-point.
The labels $\rho_t$ and $\rho_p$ are useful because
equidistant $\rho_p$ and $\rho_t$ values tend to translate to equidistant flux-surfaces
in configuration space.

\subsection{Manufactured Solution} \label{sec:manufactured}
In order to test the implementation we manufacture a solution to
Eqs.~\eqref{eq:Egyrofluid} and \eqref{eq:elliptic} of the form
\begin{align*}
\phi(R,Z,\varphi,t) &:= \sin(3\pi(R-R_0))\sin(3\pi Z)\sin(3\varphi)\sin(3\pi t)/5; \\
\psi(R,Z,\varphi,t) &:= \sin(3\pi(R-R_0))\sin(3\pi Z)\sin(3\varphi)\sin(3\pi t)/5 = \gamma_{\phi}; \\
n_e(R,Z,\varphi, t) &:= 1 + 0.5\sin(\pi(R-R_0))\sin(\pi Z)\sin(\varphi)\sin(\pi t) \\
{N_i}(R,Z,\varphi,t) &:= 1 + 0.25 \sin(\pi(R-R_0))\sin(\pi Z)\sin(\varphi)\sin(\pi t) = \gamma_{N_i} \\
u_{\parallel,e}(R,Z,\varphi, t) &:= \sin(2\pi(R-R_0))\sin(2\pi Z)\sin(2\varphi)\sin(2\pi t)/(3\sqrt{-\mu_e}) \\
U_{\parallel,i}(R,Z,\varphi, t) &:= \sin(2\pi(R-R_0))\sin(2\pi Z)\sin(2\varphi)\sin(2\pi t)/(3) \\
A_\parallel( R,Z,\varphi,t) &:= \beta( U_{\parallel,i} - u_{\parallel,e})
\end{align*}
We choose circular flux surfaces of the form
\begin{align*}
\psi_p(R,Z) :=0.5((R-R_0)^2 + Z^2),\quad
I_p(R,Z):=I_0
\end{align*}
with $R_0=10$ and $I_0=20$ and a simulation box
$[R_0-a,R_0+a]\times[-a,a]\times[0,2\pi]$.
We then insert these solutions into Eqs.~\eqref{eq:Egyrofluid} and
\eqref{eq:elliptic} and symbolically compute (with the help of Mathematica) the
remainder as source terms that are stored as C-code in the file
\mintinline{c}{manufactured.h}. We insert these source terms to the right hand
side of the corresponding equation in code and simulate for a small time-span.
By comparing the numerical solution to the manufactured one we can observe the
convergence of our numerical methods.
\begin{tcolorbox}[title=Note]
    In order to better distinguish the convergence of the DG discretized terms
    from our parallel derivative we can selectively choose to only activate
    perpendicular (including $A_\parallel$ terms) or parallel terms (those that
    involve derivatives along $\bhat$).
\end{tcolorbox}

Unfortunately, we were unable to find a closed solution for the energy integrals with the above fields.
\begin{tcolorbox}[title=Note]
    The program \mintinline{bash}{manufactured} ignores the "init" parameter
    in the input file and must be used with a circular magnetic field equilibrium
    in combination with "wall", "sheath" and "source" set as "none" or "zero".
    Further, we have the "regularization" of "order" : 1 and we have the same
    diffusion for density and velocity.
    % Check exact equations that are solved!
\end{tcolorbox}

\section{Troubleshooting} \label{sec:troubleshooting}
All previously mentioned codes can crash for various reasons. Here,
we list and describe situations, which generally may lead to program
termination
\begin{longtable}{p{6cm}p{8cm}}
\toprule
\rowcolor{gray!50}\textbf{Error condition} &  \textbf{Handling} \\ \midrule
An input file does not exist or is otherwise invalid
&
Program terminates with an error message to \mintinline{c}{std::cerr}. \mintinline{bash}{feltordiag.cu} writes an error to \mintinline{c}{std::cerr} and continues with the next input file.
    \\
An input netcdf file misses a required field
&
Program terminates with a NetCDF error message to \mintinline{c}{std::cerr}
    \\
No write permission for the output file location
&
Program terminates with an error message to \mintinline{c}{std::cerr}
    \\
An input Json file misses a key or contains a typo in a key
&
Program will exit with an error message. (The reason why we do not
silently use the default value is that the danger of wasting
valuable computing time on the cluster due to a typo is bigger than the
added convenience. We want to be sure that the program
does what the user wants).
The other programs just issue warnings
if a key is not found and use a default value
which is $0$ if not otherwise specified.
    \\
    An input Json file has an invalid value, e.g. a typo in a string value
&
Invalid values lead to termination with an error message to \mintinline{c}{std::cerr}, once and if program tries to use the value
    \\
    Number of processes in $x$, $y$ and $z$ direction does not match total number of Processes
&
Program terminates with an error message to \mintinline{c}{std::cerr}.
    \\
    $2^{s-1}$ or $c_x$ or $c_y$ does not evenly divide $N_x$ and $N_y$, where $s$ is the number of stages in the multigrid algorithm.
&
Program terminates on thrown error. Make sure the numbers add up.
    \\
    Number of processes in $x$, $y$ and $z$ direction does not evenly divide or is greater or equal $N_x/2^{s-1}$, $N_y/2^{s-1}$ and $N_z$, where $s$ is the number of stages in the multigrid algorithm.
&
Program terminates on failed assert
    \\
An MPI error occurs
&
Program crashes horribly printing cryptic error messages (stack trace) to \mintinline{c}{std::cerr}
    \\
A numerical instability occurs
&
The program terminates usually caused by a NaN exception raised. However,
the cause for the instability has to be determined inspecting the
last output in the output file.
    \\
\qquad large fieldaligned oscillations in $u_{\parallel,e}$ paired with instability in the edge of the box
&
Apply damping region
    \\
\qquad Perpendicular grid oscillations in $u_{\parallel,e}$ and $\Delta_\perp \phi$ in the damping region, symmetric in $\varphi$
&
Increase damping $alpha$, increase damping boundary, make the box larger/smaller, increasing DS refinement might help.
    \\
\qquad Spike in $u_{\parallel,e}$ shortly after simulation start
&
Increase $\nu_\perp$, increase $N_x$, $N_y$, decrease perturbation amplitude
    \\
\qquad Grid oscillations far away from the edge
&
Probably caused by the perpendicular transport that goes unstable. Increase $\nu_\perp$ and/or $N_x$, $N_y$. Increasing DS refinement might also help.
\\
\qquad Oscillations where fieldlines intersect the wall
&
Caused by boundary conditions in FCI method and necessarily underresolved toroidal direction.
Increase $N_z$, decrease $N_x$, $N_y$ or decrease $q$ value by decreasing $\mathcal P_\psi$ in geometry input file
\\
\bottomrule
\end{longtable}
%..................................................................

\bibliography{../common/references}
%..................................................................

\end{document}
